% ==============================
% Dummy Overleaf Project (HE) — Linear Algebra Book
% Exact environments: definition, example, exercise, remark,
% theorem, lemma, propositionosition, corollary, proof, solution, hint
% Numbering by chapter; solution/hint unnumbered.
% ==============================
\documentclass[12pt]{report} 

% ---- Math & friends ----
\usepackage{amsmath,amssymb,mathtools}
\usepackage{amsthm}
\allowdisplaybreaks
\usepackage{float}
\usepackage{hyperref}
\usepackage{graphicx}
\usepackage{booktabs}
\usepackage{enumitem}
\usepackage{todonotes}
\usepackage{braket}
\usepackage{array}

% ---- Hebrew + fonts ----
\usepackage{fontspec}
\usepackage{polyglossia}
\setdefaultlanguage{hebrew}
\setotherlanguage{english}
% Overleaf has Frank Ruehl CLM in TeX Live
\newfontfamily\hebrewfont[Script=Hebrew]{Frank Ruehl CLM}
\newfontfamily\hebrewfontsf[Script=Hebrew]{Frank Ruehl CLM}
\newfontfamily\hebrewfonttt[Script=Hebrew]{Frank Ruehl CLM}


% ---- Theorem-like environments (by chapter) ----
\theoremstyle{plain}
\newtheorem{theorem}{משפט}[chapter]
\newtheorem{lemma}[theorem]{למה}
\newtheorem{proposition}[theorem]{טענה}
\newtheorem{corollary}[theorem]{מסקנה}

\theoremstyle{definition}
\newtheorem{definition}[theorem]{הגדרה}
\newtheorem{example}{דוגמה}[chapter]
\newtheorem{exercise}{תרגיל}[chapter]
\newtheorem{remark}{הערה}[chapter]

% Unnumbered for website folding
\newtheorem*{solution}{פתרון}
\newtheorem*{hint}{רמז}

% ---- Macros (optional) ----
\newcommand{\R}{\mathbb{R}}
\newcommand{\N}{\mathbb{N}}
\newcommand{\Z}{\mathbb{Z}}
\newcommand{\Q}{\mathbb{Q}}
\newcommand{\C}{\mathbb{C}}
\newcommand{\F}{\mathbb{F}}
\DeclareMathOperator{\rank}{rank}
\newcommand{\M}{\mathbb{M}}
\newcommand{\0}{\mathbf{0}}


\title{ספר אלגברה לינארית — דגם אוברליף}
\author{צוות הקורס}
\date{\today}

\setlength {\marginparwidth }{2cm} 

\begin{document}
\maketitle

\chapter{מטריצות}\label{ch:matrices}

בפרק הקודם ראינו מטריצות מקדמים בהקשר של ממ"ל (מצומצמות או מורחבות). אבל הרעיון של מטריצה הוא כללי יותר, ולא תמיד חייב להיות קשר לממ"ל. מטריצה אינה אלא טבלה של מספרים, וטבלה כזו יכולה לייצג כל מיני דברים.

\begin{example}\label{ex:bus}
נניח שאנחנו רוצים לחשב את מספר האוטובוסים (של חברה מסוימת) שעוברים בין אילת, באר שבע ות"א ביום ראשון. לצורך הדוגמה, לכל אוטובוס אין תחנות ביניים - רק תחנת יציאה ותחנת הגעה. אז אפשר להכין מטריצה מסדר 
$3\times 3$
שבה כל שורה תציין תחנת יציאה, וכל עמודה תציין תחנת הגעה. נשתמש בסימונים 
$E, B, T$
לת"א, באר שבע ואילת בהתאמה כדי שיהיה ברור איזו עמודה/שורה מתאימה לאיזו עיר (סימונים אלה אינם איברים במטריצה). אז לפי הנתונים מקבלים מטריצה כזו:


$$
\left(
\begin{array}{c|ccc}
 & E & B & T \\ \hline
E & 0 & 10 & 5 \\
B & 8 & 0 & 12 \\
T & 4 & 6 & 0
\end{array}
\right)
$$

משמעות המטריצה היא שאין אוטובוס מאף תחנה לעצמה. יש עשרה אוטובוסים שיוצאים מאילת לבאר שבע, ויש שמונה אוטובוסים שעושים את המסלול ההפוך. וכן הלאה. נוותר על הסימונים של הערים ונכתוב את המטריצה בצורה הרגילה, ונסמנה ע"י $S$:

$$
S=\left(
\begin{array}{ccc}
0 & 10 & 5 \\
8 & 0 & 12 \\
4 & 6 & 0
\end{array}
\right)
$$

זו המטריצה עבור הנתונים של יום ראשון. עכשיו נניח שיש נתונים שונים ליום שני, ונכין מטריצה חדשה עבורם:

$$
M=\left(
\begin{array}{ccc}
0 & 12 & 7 \\
11 & 0 & 10 \\
3 & 8 & 0
\end{array}
\right)
$$

כעת אפשר לשאול כמה אוטובוסים יש לכל מסלול בשני הימים יחד. אז צריך לעשות חיבור בכל רכיב של המטריצה (שמתאים למסלול עם נקודת יציאה ונקודת הגעה) לחוד. אפשר לכתוב מטריצה חדשה שאיבריה יהיו סכומי האיברים של המטריצות
$S,M$
לפי כל רכיב לחוד. זה מתבקש לקרוא למטריצה החדשה 
$S+M$
בדומה לחיבור וקטורי. כך נקבל:

$$
S+M=\left(
\begin{array}{ccc}
0+0 & 10+12 & 5+7 \\
8+11 & 0+0 & 12+10 \\
4+3 & 6+8 & 0+0
\end{array}
\right)
=\left(
\begin{array}{ccc}
0+0 & 22 & 12 \\
19 & 0 & 22 \\
7 & 14 & 0
\end{array}
\right)
$$

בפרט, משמעות המטריצה לעיל היא שאין אוטובוס מאף תחנה לעצמה באף יום. יש
$22$
אוטובסים שיוצאים מאילת לבאר שבע ביומיים של ראשון ושני, ו-
$19$
שעושים את המסלול ההפוך ביומיים האלה. וכן הלאה.

\end{example}

\begin{definition}
\adddeftolist{מרחב המטריצות $\M_{m\times n}(\F)$}
נגדיר את קבוצת כל המטריצות מסדר 
$m\times n$
מעל
$\F$
(כאשר
$\F=\R$
או
$\F=\C$
)
ע"י 

$$
.\M_{m\times n}(\F)=\Set{\begin{pmatrix}a_{11} & a_{12} & \cdots & a_{1n}\\
a_{21} & a_{22} & \cdots & a_{2n}\\
\vdots & \vdots & \ddots & \vdots\\
a_{m1} & a_{m2} & \cdots & a_{mn}
\end{pmatrix}|\forall 1\leq i\leq m \quad \forall 1\leq j\leq n \quad a_{ij}\in \F}
$$

\end{definition}

\begin{remark}
הסימן 
$\forall$
פירושו "לכל". אז התנאי המתמטי בהגדרת 
$\M_{m\times n}(\F)$
אומר שכל איבר 
$a_{ij}$
במטריצה שייך לקבוצה 
$\F$,
כי עוברים על כל האיברים בכל השורות ובכל העמודות.
\end{remark}

\begin{example}\
\begin{enumerate}[label=\alph*.]

\item
מתקיים
$$
.A=\begin{pmatrix} 10 & 9 & 1 & -7\\
1 & 2 & 3 & 4\\
0 & 4 & -3 & -5
\end{pmatrix}\in\M_{3\times 4}(\R)
$$


\item
מתקיים
$$
.B=\begin{pmatrix} i & 1 \\
1 & 1+i\\
0 & 2i 
\end{pmatrix}\in\M_{3\times 2}(\C)
$$

\begin{remark}
לכל
$m,n\in\N$
מתקיים
$\M_{m\times n}(\R)\subseteq \M_{m\times n}(\C)$ 
כי
$\R\subseteq\C$. 
אז למעשה כל המטריצות בקורס הן מעל $\C$,
אבל אם לדייק רובן (לא כולן) יהיו מעל $\R$.
השימוש בסימון $\F$ 
נועד לרוב הטענות, שנכונות גם מעל $\R$
וגם מעל $\C$.
בתחום של אלגברה לינארית אפשר גם לעסוק בקבוצות מספרים נוספות שנקראות שדות
,(fields)
ומכאן הסיבה לאות 
$\F$.
אבל בקורס שלנו נסתפק ב-
$\R,\C$
(השדות החשובים ביותר).
\end{remark}


\end{enumerate}
\end{example}

\section{פעולות על מטריצות}
\subsection{חיבור וכפל בסקלר}

ראינו בדוגמה
\ref{ex:bus}
שיש טעם בהגדרת פעולת חיבור בין מטריצות, בדיוק כמו חיבור וקטורי (עושים חיבור בכל קוארדינטה/רכיב לחוד). קודם כל נכליל חיבור וקטורי וכפל בסקלר (שהגדרנו בפרק
\ref{ch:geometry}
במישור ובמרחב) 
ל-
$\F^n$.

\begin{definition}
\adddeftolist{חיבור וקטורי וכפל בסקלר במרחב $\F^n$}
עבור וקטורים
$ה=(a_1,...,a_n)\in\F^n,w=(b_1,...,b_n)\in\F^n$
וסקלר
$c\in\F$
נגדיר את הפעולות הבאות:
\begin{enumerate}[label=\alph*.]
\item 
פעולת החיבור:
$$
v+w=(a_1+b_1,...,a_n+b_n)
$$


\item
פעולת כפל בסקלר:
$$
c \cdot v=(c\cdot a_1,\dots,c\cdot a_n)
$$


\end{enumerate}
\end{definition}

\begin{definition}
\adddeftolist{חיבור מטריצות וכפל בסקלר}
עבור
$A,B\in\M_{m\times n}(\F)$
וסקלר
$c\in\F$
נגדיר את הפעולות הבאות:

\begin{enumerate}[label=\alph*.]
\item 
פעולת החיבור:
$$
\begin{aligned}
A+B&=\begin{pmatrix}a_{11} & \cdots & a_{1n}\\
a_{21} & \cdots & a_{2n}\\
\vdots  & \ddots & \vdots\\
a_{m1}  & \cdots & a_{mn}
\end{pmatrix}+\begin{pmatrix}b_{11} & \cdots & b_{1n}\\
b_{21} & \cdots & b_{2n}\\
\vdots & \ddots & \vdots\\
b_{m1} & \cdots & b_{mn}
\end{pmatrix}\\
&=\begin{pmatrix}a_{11}+b_{11} & \cdots & a_{1n}+b_{1n} \\ a_{21}+b_{21} & \cdots & a_{2n}+b_{2n}\\
\vdots & \ddots & \vdots\\
a_{m1}+b_{m1} & \cdots & a_{mn}+b_{mn}
\end{pmatrix}
\end{aligned}
$$


\item
פעולת כפל בסקלר:
$$
c A=\begin{pmatrix}ca_{11} & ca_{12} & \cdots & ca_{1n}\\
ca_{21} & ca_{22} & \cdots & ca_{2n}\\
\vdots & \vdots & \ddots & \vdots\\
ca_{m1} & ca_{m2} & \cdots & ca_{mn}
\end{pmatrix}
$$

בפרט, נגדיר
$-A=(-1)\cdot A$ 
ובהתאמה 
$B-A=B+(-1)\cdot A$.


\end{enumerate}
\end{definition}

\begin{example}
    \begin{enumerate}[label=\alph*.]
        \item
חיבור של מטריצות מסדר
$2\times 2$
מעל $\R$:
$$
\begin{aligned}
&\begin{pmatrix} 3 & 7 \\
5 & -2 
\end{pmatrix}+\begin{pmatrix} -10 & 5 \\
-4 & -3 
\end{pmatrix}\\
&=\begin{pmatrix} 3-10 & 7+5 \\
5-4 & -2-3 
\end{pmatrix}=\begin{pmatrix} -7 & 12 \\
1 & -5 
\end{pmatrix}
\end{aligned}
$$


\item
חיבור של מטריצות מסדר
$3\times 4$
מעל $\R$:
$$
\begin{aligned}
&\begin{pmatrix} 10 & 9 & 1 & -7\\
1 & 2 & 3 & 4\\
0 & 4 & -3 & -5
\end{pmatrix}+\begin{pmatrix} 3 & -2 & 0 & 6\\
5 & -2 & 4 & -4\\
10 & -6 & 3 & 1
\end{pmatrix}\\
&=\begin{pmatrix} 10+3 & 9-2 & 1+0 & -7+6\\
1+5 & 2-2 & 3+4 & 4-4\\
0+10 & 4-6 & -3+3 & -5+1
\end{pmatrix}\\
&=\begin{pmatrix} 13 & 7 & 1 & -1\\
6 & 0 & 7 & 0\\
10 & -2 & 0 & -4
\end{pmatrix}
\end{aligned}
$$


\item
חיבור של מטריצות מסדר
$2\times 2$
מעל $\C$:
אם

$$ A = \begin{pmatrix} 1 & i \\
2-3i & 3-2i 
\end{pmatrix},\quad  B = \begin{pmatrix} 7-5i & 2+7i \\
-6 & 3+10i 
\end{pmatrix}$$

אז
$$
 \begin{aligned}
A+B&=\begin{pmatrix} 1+7-5i & i+2+7i \\
2-3i-6 & 3-2i+3+10i 
\end{pmatrix}=\begin{pmatrix} 8-5i & 2+8i \\
-4-3i & 6+8i 
\end{pmatrix}
\end{aligned}
$$

    \end{enumerate}
\end{example}

\begin{definition}\label{remark:entry-formula}
יש דרך מקוצרת לכתיבת חיבור. במקום לכתוב את החיבור בכל רכיב, נסתכל על האיבר הכללי של כל מטריצה. נסמן את האיבר בשורה ה-
$i$
ובעמודה ה-
$j$
ע"י
$(P)_{ij}$
לכל
$P\in\M_{m\times n}(\F)$. 
אז למעשה הגדרנו
לכל
$1\leq i \leq m$
ולכל
$1\leq j\leq n$:
$$(A+B)_{ij}=(A)_{ij}+(B)_{ij}$$

$$(c A)_{ij}=c(A)_{ij}$$

\end{definition}

למשל, בדוגמה האחרונה היה לנו
$(A)_{12} = i, (B)_{21} = -6, (A+B)_{22}=6+8i$.

\begin{remark}
שימו לב שהחיבור מוגדר רק למטריצות מאותו הסדר (כנ"ל לוקטורים - רק אם הם שווי אורך). גם הסדר של מטריצת הסכום שווה לסדר המשותף של המטריצות.
\end{remark}


\begin{example}
    החיבור הבא אינו מוגדר:
$$
\begin{pmatrix} 3 & 7 \\
5 & -2 
\end{pmatrix}+\begin{pmatrix} -10 & 5 &0\\
-4 & -3 &0
\end{pmatrix}
$$

המטריצה השמאלית היא מסדר
$2\times 2$
והימנית היא מסדר
$2\times 3$.

\end{example}
\begin{example}[דוגמאות נוספות] \ 
\begin{enumerate}[label=\alph*.]


\item 

כפל בסקלר של מטריצה מסדר
$3\times 3$
מעל $\C$:
$$
\begin{aligned}
&i\begin{pmatrix} 1 & i & 1+i \\ 
1-i & 2i & 3 \\
2+3i & 3-2i & 5
\end{pmatrix}=\begin{pmatrix} i\cdot1 & i\cdot i & i(1+i) \\ 
i(1-i) & i\cdot 2i & i\cdot3 \\
i(2+3i) & i(3-2i) & i\cdot5
\end{pmatrix}\\
&=\begin{pmatrix}i & -1 & -1+i \\ 
1+i & -2 & 3i \\
-3+2i & 2+3i & 5i
\end{pmatrix}
\end{aligned}
$$


\item 
שילוב של כפל בסקלר וחיבור מטריצות מסדר
$2\times 2$
מעל $\R$:
$$
\begin{aligned}
&2 \begin{pmatrix} 11 & 2 \\
-2 & 4 
\end{pmatrix}+3\begin{pmatrix} -4 & 5 \\
6 & -3 
\end{pmatrix}=\begin{pmatrix} 22 & 4 \\
-4 & 8 
\end{pmatrix}+\begin{pmatrix} -12 & 15 \\
18 & -9 
\end{pmatrix}\\
&=\begin{pmatrix} 22-12 & 4+15 \\
-4+18 & 8-9 
\end{pmatrix}
\end{aligned}
$$


\end{enumerate}
\end{example}

\begin{definition}
\adddeftolist{מטריצת האפס $0_{m\times n}$}
לכל 
$m,n\in\N$
נגדיר את מטריצת האפס מסדר
$m\times n$
כמטריצה שכל איבריה הם
$0$:
$$
.\0_{m\times n}=\begin{pmatrix}0 &0 &\cdots &0 \\
0 &0 &\cdots &0 \\
\vdots & \vdots &\ddots &\vdots\\
0 &0 &\cdots &0
\end{pmatrix}
$$
\end{definition}

\begin{proposition}
\addpropositiontolist{תכונות של פעולות על מטריצות}
יהיו 
$A,B,C\in\M_{m\times n}(\F)$
ו-
$\alpha,\beta\in\F$.
אז מתקיימים הכללים הבאים:

\begin{enumerate}[label=\alph*.]
\item
חוק החילוף (קומוטטיביות): 
$$A+B=B+A$$


\item 
חוק הקיבוץ (אסוציאטיביות):
$$(A+B)+C=A+(B+C)$$


\item 
חוק הקיבוץ לכפל בסקלר:
$$(\alpha\beta)A=\alpha(\beta A)$$


\item 
חוק הפילוג (דיסטריביוטיביות):
$$\alpha(A+B)=\alpha A+\alpha B$$


\item
ניטרליות היחידה:
$$1\cdot A=A$$


\item
$$0\cdot A=\0_{m\times n}$$


\item 
ניטרליות מטריצת האפס:
$$A+\0_{m\times n}=A$$


\item 
$$A-A=0_{m\times n}$$


\end{enumerate}
\end{proposition}

\begin{proof}
כל התכונות עוברות בירושה מחיבור וכפל ב-
$\F$
(מדובר בתכונות מוכרות של 
$\R$
וראינו בפרק 
\ref{ch:complex}
שהן עוברות בירושה ל-
$\C$).
נוכיח למשל את חוק הפילוג: נתמקד באיבר הכללי של כל מטריצה בשוויון שצריך להוכיח.
לכל
$1\leq i\leq m, 1\leq j\leq n$
מתקיים

$$
.(\alpha(A+B))_{ij}=\alpha(A+B)_{ij}=\alpha((A)_{ij}+(B)_{ij})=\alpha (A)_{ij}+\alpha (B)_{ij}
$$
שימו לב שהמעבר האחרון נכון לפי חוק הפילוג ב-$\F$.
לכן נובע כי
$$.\alpha(A+B)=\alpha A+\beta A$$

נוכיח גם את ניטרליות מטריצת האפס:
$$
(A+\0_{m\times n})_{ij}=(A)_{ij}+(\0_{m\times n})_{ij}=(A)_{ij}+0=(A)_{ij}
$$
לכל
$1\leq i\leq m, 1\leq j\leq n$,
ולכן
$$,A+\0_{m\times n}=A$$

ההוכחות של שאר הסעיפים דומות, כאשר הרעיון הכללי הוא שימוש בהערה
\ref{remark:entry-formula}
והתכונות המתאימות של חיבור/כפל ב-
$\F$.

\end{proof}

\begin{exercise}
הוכיחו את חוק הקיבוץ לכפל בסקלר.
\end{exercise}

\begin{solution}
עבור חוק הקיבוץ: לכל 
$1\leq i\leq m, 1\leq j\leq n$
מתקיים 

$$
.((A+B)+C)_{ij}=(A+B)_{ij}+(C)_{ij}=((A)_{ij}+(B)_{ij})+(C)_{ij}
$$

באופן דומה, מתקיים
$$
.(A+(B+C))_{ij}=A_{ij}+(B+C)_{ij}=(A)_{ij}+((B)_{ij}+(C)_{ij})
$$

שני הביטויים באגף ימין של שתי המשוואות שווים לפי חוק הקיבוץ ב-
$\F$.
לכן
$$.(A+B)+C=A+(B+C)$$
\end{solution}

\subsection{כפל מטריצה בוקטור עמודה}

בפרק הקודם דיברנו על ממ"ל מהצורה
$Ax=b$
כאשר
$$
A=\left( \begin{array}{@{}cccc}
	a_{11} & \cdots & a_{1n}  \\
	a_{21} & \cdots & a_{2n} \\
	\vdots & \ddots & \vdots  \\
	a_{m1}  & \cdots & a_{mn} 
\end{array} \right)
$$
היא מטריצת המקדמים המצומצמת, 
$x=\begin{pmatrix}x_1\\
x_2\\
\vdots\\
x_n
\end{pmatrix}$
הוא וקטור המשתנים ו-
$b=\begin{pmatrix}b_1\\
b_2\\
\vdots\\
b_m
\end{pmatrix}$
הוא וקטור הקבועים. 
נרצה להגדיר באופן מדויק את פעולת הכפל 
$Ax$
של מטריצה 
$A\in\M_{m\times n}(\F)$
בוקטור עמודה
$x\in\F^n$
שמספר שורותיו שווה למספר העמודות של המטריצה.

לפני כן, נגדיר את סימון 
$\Sigma$
לסכימה (שאולי כבר מוכר).

\begin{definition}
\adddeftolist{סכימה עם $\Sigma$}
בהינתן
$a_1,a_2,...,a_n\in\F$
נסמן את סכומם בעזרת כתיבה עם
$\Sigma$
ואינדקס
$l$.

$$
.\sum_{l=1}^n a_l=a_1+a_2+a_3+\cdots+a_n
$$

האינדקס
$l$
רץ על פני הקבוצה
$\Set{1,2,3,\cdots,n}$
וכל הצבה תורמת איבר לסכום.

\end{definition}

\begin{example}
במקרה
$a_l=l$
נקבל
$$
\sum_{l=1}^{5} l=1+2+3+4+5=15
$$
\end{example}

\begin{definition}
\adddeftolist{כפל מטריצה בוקטור עמודה}
בהינתן
$A\in\M_{m\times n}(\F)$
ו-
$$
,x=\begin{pmatrix}
x_1\\
x_2\\
\vdots \\
x_n
\end{pmatrix}\in\F^n
$$
נגדיר את המכפלה
$$
Ax=\begin{pmatrix}
y_1\\
y_2\\
\vdots \\
y_m
\end{pmatrix}\in\F^m
$$
עם קוארדינטות הנתונות ע"י
$$
y_i=\sum_{j=1}^n (A)_{ij}x_j
$$
לכל
$1\leq i\leq m$.
נסמן
$y_i=(Ax)_i$ 
עבור הקוארדינטות.
\end{definition}

\begin{remark}
לפי הגדרה זו, לכל 
$1\leq i\leq m$
מתקיים
$$
.(Ax)_i=b_i \iff \sum_{j=1}^n (A)_{ij}x_j=b_i
$$

לכן, השוויון הוקטורי
$Ax=b$
אכן שקול לממ"ל.
\end{remark}

\begin{remark}
למעשה יש זיהוי בין
$\M_{1\times n}(\F)$
לבין 
$\F^n$,
לפי הגדרה
\ref{def:F^n} 
שהתייחסה לוקטורי שורה. אבל נרשה לעצמנו להשתמש בסימון 
$\F^n$
גם עבור וקטורי עמודה באורך
$n$,
כלומר עבור הקבוצה
$\M_{n\times 1}(\F)$. 
לפעמים ההבדל בין וקטור שורה לוקטור עמודה הוא אסתטי בלבד,
אבל כאשר מכפילים מטריצה בוקטור בהחלט יש חשיבות לסוג הוקטור (הוא צריך להיות וקטור עמודה). אז ככלל נזכור לשים לב להבדל כאשר מופיעה מכפלה של מטריצה בוקטור, ובשאר המקרים אפשר לכתוב את הוקטור כשורה או עמודה לפי טעם אישי.
\end{remark}

\begin{example}\
\begin{enumerate}[label=\alph*.]
\item
ניקח
$$
A = \begin{pmatrix}
2 & -1 \\
3 & 4
\end{pmatrix}, \quad 
x = \begin{pmatrix}
1 \\
5
\end{pmatrix}
$$
ונקבל
$$
.\begin{pmatrix}
2 & -1 \\
3 & 4
\end{pmatrix}
\begin{pmatrix}
1 \\
5
\end{pmatrix}
=
\begin{pmatrix}
2\cdot1 + (-1)5 \\
3\cdot1 + 4\cdot5
\end{pmatrix}
=
\begin{pmatrix}
-3 \\
23
\end{pmatrix}
$$


\item 
$$
\begin{pmatrix}a_1 &b_1 &c_1\end{pmatrix}\begin{pmatrix}a_2 \\ b_2 \\ c_2\end{pmatrix}=a_1a_2+b_1b_2+c_1c_2
$$

לכל 
$a_1,b_1,c_1,a_2,b_2,c_2\in\F$. 
במקרה של
$\F=\R$
מדובר על המכפלה הסקלרית של
$$
.\begin{pmatrix}a_1 \\ b_1 \\ c_1\end{pmatrix},\begin{pmatrix}a_2 \\ b_2 \\ c_2\end{pmatrix}\in\R^3
$$


\item 
$$
\begin{pmatrix}
1 & 0 & 2 \\
-2 & 3 & 1 \\
4 & -1 & 0
\end{pmatrix}
\begin{pmatrix}
3 \\
2 \\
-1
\end{pmatrix}
=
\begin{pmatrix}
1\cdot3 + 0\cdot2 + 2(-1) \\
-2\cdot3 + 3\cdot 2 + 1(-1) \\
4\cdot3 + (-1)2 + 0(-1)
\end{pmatrix}
=
\begin{pmatrix}
1 \\
-1 \\
10
\end{pmatrix}
$$

\item
$$
\begin{pmatrix}
1+i & 2 \\
3i & 4-i
\end{pmatrix}\begin{pmatrix}
2 \\
1-i
\end{pmatrix}=\begin{pmatrix}
(1+i)2 + 2(1-i) \\
(3i)2 + (4-i)(1-i)
\end{pmatrix}
=
\begin{pmatrix}
4 \\
3 + i
\end{pmatrix}
$$


\end{enumerate}
\end{example}

\begin{exercise}
חשבו את 
$Av$
עבור
$$
.A = \begin{pmatrix}
1 & i & 2 \\
-1 & 3i & 0 \\
2+i & 1 & -i
\end{pmatrix}, \quad 
v = \begin{pmatrix}
1+i \\
2 \\
-1
\end{pmatrix}
$$
\end{exercise}

\begin{solution}
$$
Ax=\begin{pmatrix}
1(1+i) + i\cdot 2 + 2(-1) \\
(-1)(1+i) + 3i\cdot2 + 0(-1) \\
(2+i)(1+i) + 1\cdot2 + (-i)(-1)
\end{pmatrix}
=
\begin{pmatrix}
1+i + 2i - 2 \\
-1 - i + 6i \\
1+3i+2+i
\end{pmatrix}
=
\begin{pmatrix}
-1 + 3i \\
-1 + 5i \\
3+4i
\end{pmatrix}
$$
\end{solution}

\begin{proposition}\label{prop:matrix-vector}
\addpropositiontolist{תכונות של כפל מטריצה בוקטור עמודה}
יהיו
$A,B\in\M_{m\times n}(\F)$,
$v,w\in\F^n$
ו-
$\alpha\in\F$.
\begin{enumerate}[label=\alph*.]
\item
$$
A(v+w)=Av+Aw
$$


\item 
$$
(A+B)v=Av+Bv
$$


\item 
$$
A(\alpha v)=\alpha(Av)
$$


\end{enumerate}
\end{proposition}

\begin{proof}
בכל סעיף נראה שוויון בין הקוארדינטות ה-
$i$
של שני האגפים, לכל
$1\leq i\leq m$. 
ראשית נסמן
$$
.v=\begin{pmatrix}
 v_1 \\
 v_2 \\
 \vdots \\
 v_n
 \end{pmatrix}, \quad 
 w=\begin{pmatrix}
 w_1 \\
 w_2 \\
 \vdots \\
 w_n 
 \end{pmatrix}
$$

\begin{enumerate}[label=\alph*.]
\item 
$$
(A(v+w))_i=\sum_{j=1}^n (A)_{ij}(v+w)_j=\sum_{j=1}^n (A)_{ij}v_j+\sum_{j=1}^n (A)_{ij}w_j=(Av)_i+(Aw)_i
$$
בשביל השוויון השני השתמשנו בחוק הפילוג ב-
$\F$, 
וגם פיצלנו את הסכום לפי חוק החילוף (אין חשיבות לסדר הסכימה).


\item 
$$
\begin{aligned}
((A+B)v)_i&=\sum_{j=1}^n ((A)_{ij}+(B)_{ij})v_j=\sum_{j=1}^n (A)_{ij}v_j+\sum_{j=1}^n (B)_{ij}v_j\\
&=(Av)_i+(Bv)_i=(Av+Bv)_i
\end{aligned}
$$


\item 
$$
(A(\alpha v))_i=\sum_{j=1}^n (A)_{ij}(\alpha v_j)=\alpha\sum_{j=1}^n (A)_{ij}v_j=\alpha(Av)_i=(\alpha Av)_i
$$
בשוויון השני השתמשנו בחוק הפילוג ב-
$\F$,
לפיו אפשר להוציא גורם משותף מחוץ לסכימה.

\end{enumerate}
\end{proof}

בהינתן מטריצה
$A$,
אפשר לחשוב על הכפל
$Av$
כפעולה על
$v$.
זהו רעיון שנחזור אליו בהמשך הקורס, אבל כבר ניתן להתרשם מהפעולה הזו באופן גיאומטרי, למשל עבור
$$
.A=\begin{pmatrix}
1 & 1 \\
0 & -2
\end{pmatrix}
$$

אפשר לצייר רשת במרחב המקור (עבור
$v$)
בהתאם לוקטורים
$$
\begin{pmatrix}1 \\
0 
\end{pmatrix}, \begin{pmatrix}0 \\
1
\end{pmatrix}
$$
ולראות איזו רשת נוצרת במרחב התמונה (עבור
$Av$):

\begin{figure}[H]
    \centering
    \includegraphics[width=0.7\textwidth]{images/MatrixGrid.png}
    \caption{הכפל במטריצה משפיע על שני הוקטורים באופן שונה, אבל כל ריבוע במרחב המקור הופך למקבילית במרחב התמונה}
    \label{fig:MatrixGrid}
\end{figure}

במקום רשת של ריבועים, אפשר להסתכל על מעגלים סביב הראשית ולראות איזה מסלול מתאים לכל מעגל במרחב התמונה:

\begin{figure}[H]
    \centering
    \includegraphics[width=0.7\textwidth]{images/MatrixCircle.png}
    \caption{כל מעגל במרחב המקור נשלח לאליפסה במרחב התמונה, ומתיחה של המעגל מתאימה למתיחה של האליפסה}
    \label{fig:MatrixCircle}
\end{figure}

\begin{proposition}
\addpropositiontolist{תכונות של קבוצת פתרונות לממ"ל הומוגנית}
תהי
$Ax=\underline{0}$
ממ"ל הומוגנית.
\begin{enumerate}[label=\alph*.]
\item
קבוצת הפתרונות סגורה לחיבור, כלומר לכל שני פתרונות
$v,w\in\F^n$
גם
$v+w$
הוא פתרון.


\item
קבוצת הפתרונות סגורה לכפל בסקלר, כלומר לכל פתרון
$v\in\F^n$
ולכל
$\alpha\in\F$
גם
$\alpha v$
הוא פתרון.


\end{enumerate}
\end{proposition}

\begin{proof}
יהיו 
$v,w\in\F^n$
פתרונות של הממ"ל.
\begin{enumerate}[label=\alph*.]
\item
 לפי טענה
\ref{prop:matrix-vector}
מתקיים
$$
.A(v+w)=Av+Aw=\underline{0}+\underline{0}=\underline{0}
$$
כלומר
$v+w$
הוא פתרון של הממ"ל.



\item 
שוב, לפי טענה
\ref{prop:matrix-vector}
לכל 
$\alpha\in\F$
מתקיים

$$
.A(\alpha v)=\alpha Av=\alpha\cdot 0=0
$$
כלומר 
$\alpha v$
הוא פתרון של הממ"ל.


\end{enumerate}
\end{proof}

\begin{proposition}
\addpropositiontolist{הקשר בין קבוצות הפתרונות של ממ"ל לא הומוגנית והומוגנית}
תהי
$Ax=b$
ממ"ל לא הומוגנית. נניח שקיים לה פתרון (לא בהכרח יחיד)
$v_0\in\F^n$.
אז מתקיים הקשר הבא בין קבוצת הפתרונות של הממ"ל הלא הומוגנית לקבוצת הפתרונות של הממ"ל ההומוגנית
$Ax=\underline{0}$:

$$
.\Set{x\in\F^n|Ax=b}=\Set{v_0+y|y\in\F^n,Ay=\underline{0}}
$$
\end{proposition}

\begin{remark}
הרעיון הזה כבר מוכר לנו מההצגה הפרמטרית של ישר/מישור. למשל, ראינו שלישר במישור יש הצגה פרמטרית מהצורה
$$
\Set{\begin{pmatrix}
x_0 \\
y_0 
\end{pmatrix}+t\begin{pmatrix}
\alpha \\
\beta
\end{pmatrix}|t\in\R}
$$
כאשר
$\begin{pmatrix}
\alpha \\
\beta
\end{pmatrix}\in\R^2$
הוא וקטור כיוון (בפרק 
\ref{ch:geometry}
השתמשנו בוקטורי שורה, אבל בהקשר של כפל צריך לקחת וקטורי עמודה). אם הישר מתאר קבוצת פתרונות של ממ"ל
$\begin{pmatrix}a_1 &a_2\end{pmatrix}\begin{pmatrix}
x \\
y 
\end{pmatrix}=b$,
אז וקטור הכיוון הוא פתרון של הממ"ל ההומוגנית. כלומר מתקיים
$\begin{pmatrix}a_1 &a_2\end{pmatrix}\begin{pmatrix}
\alpha \\
\beta
\end{pmatrix}=0$,
או באופן שקול
$a_1\alpha+a_2\beta=0$.
יותר מכך, קבוצת הפתרונות של הממ"ל ההומוגנית היא קבוצת הכפולות בסקלר של וקטור הכיוון.
\end{remark}

נוכיח את הטענה:

\begin{proof}
נוכיח שוויון בין שתי הקבוצות ע"י הכלה דו-כיוונית.

הכיוון
$\supseteq$:
יהי 
$y\in\F^n$
כך ש-
$Ay=\underline{0}$.
אז מתקיים
$$
,A(v_0+y)=Av_0+Ay=b+0=b
$$
ולכן
$v_0+y\in\Set{x\in\F^n|Ax=b}$.

הכיוון
$\subseteq$:
יהי 
$x\in\F^n$
כך ש-
$Ax=b$.
אז מתקיים
$x=v_0+y$
עבור
$y=x-v_0$.
נבדוק ש-$y$
הוא פתרון של הממ"ל ההומוגנית:

$$
.Ay=A(x-v_0)=Ax-Av_0=b-b=\underline{0}
$$

לכן

$$
.x=v_0+y\in\Set{x_0+y|y\in\F^n,Ay=\underline{0}}
$$

\end{proof}

\subsection{כפל מטריצות}

\begin{definition}
\adddeftolist{כפל מטריצות}
בהינתן
$A\in\M_{m\times k}(\F),B\in\M_{k\times n}(\F)$
נגדיר את מטריצת הכפל
$AB$
ע"י
$$
(AB)_{ij}=\sum_{l=1}^k (A)_{il}(B)_{lj}=(A)_{i1}(B)_{1j}+(A)_{i2}(B)_{2j}+...+(A)_{ik}(B)_{kj}
$$
לכל
$1\leq i\leq m, 1\leq j\leq n$.
\end{definition}

\begin{remark}
שימו לב שכפל בין מטריצות מוגדר רק כאשר מספר העמודות של המטריצה השמאלית שווה למספר השורות של המטריצה הימנית, וסימנו את מספר זה ע"י
$k$. 
זה אומר שבהרבה מקרים פעולת הכפל אינה מוגדרת, למשל כאשר
$A,B$
שתיהן מסדר
$2\times 3$.
\end{remark}

\begin{example}\
\begin{enumerate}[label=\alph*.]
\item 
עבור
$
A=\begin{pmatrix}
2 & 3 \\
1 & 4
\end{pmatrix},B=
\begin{pmatrix}
5 & 6 \\
7 & 8
\end{pmatrix}
$
נקבל
$$
AB=\begin{pmatrix}
2\cdot5 + 3\cdot7 & 2\cdot6 + 3\cdot8 \\
1\cdot5 + 4\cdot7 & 1\cdot6 + 4\cdot8
\end{pmatrix}
=
\begin{pmatrix}
31 & 36 \\
33 & 38
\end{pmatrix}
$$

לעומת זאת, אם נחשב את המכפלה בסדר הפוך נקבל

$$
BA=\begin{pmatrix}
5 & 6 \\
7 & 8
\end{pmatrix}\begin{pmatrix}
2 & 3 \\
1 & 4
\end{pmatrix}
=
\begin{pmatrix}
5\cdot2+6\cdot1 & 5\cdot3+6\cdot4 \\
7\cdot2+8\cdot1 & 7\cdot3+8\cdot4
\end{pmatrix}
=
\begin{pmatrix}
16 & 39 \\
22 & 53
\end{pmatrix}
$$
אז במקרה זה
$AB\neq BA$.


\item 
עבור
$
C=\begin{pmatrix}
1 & 2 & 3 \\
4 & 5 & 6
\end{pmatrix},\quad D=
\begin{pmatrix}
1 & 0 & 2 \\
-1 & 3 & 1 \\
0 & 4 & 5
\end{pmatrix}
$
נקבל
$$
\begin{aligned}
CD&=
\begin{pmatrix}
1\cdot1 + 2(-1) + 3\cdot0 & 1\cdot0 + 2\cdot3 + 3\cdot4 & 1\cdot2 + 2\cdot 1 + 3\cdot5 \\
4\cdot1 + 5(-1) + 6\cdot0 & 4\cdot0 + 5\cdot3 + 6\cdot4 & 4\cdot2 + 5\cdot1 + 6\cdot5
\end{pmatrix}\\
&=
\begin{pmatrix}
-1 & 18 & 19 \\
-1 & 39 & 43
\end{pmatrix}
\end{aligned}
$$


\item
עבור 
$E=\begin{pmatrix}
1+i & 2 \\
3 & 4-i
\end{pmatrix}, \quad
F=\begin{pmatrix}
2 & i \\
1-i & 3
\end{pmatrix}$
נקבל
$$
\begin{aligned}
EF&=\begin{pmatrix}
(1+i)2 + 2(1-i) & (1+i)i + 2\cdot3 \\
3\cdot2 + (4-i)(1-i) & 3\cdot i + (4-i)3
\end{pmatrix}\\
&=
\begin{pmatrix}
2+2i + 2 - 2i & i - 1 + 6 \\
6 + 4 - 4i - i - 1 & 3i + 12 - 3i
\end{pmatrix}=
\begin{pmatrix}
4 & 5 + i \\
9 - 5i & 12
\end{pmatrix}
\end{aligned}
$$


\item 

המכפלה
$$
\begin{pmatrix}
2 &3 &1 \\
0 &-1 &-2
\end{pmatrix}
\begin{pmatrix}
10 &-5 \\
3 &-1 
\end{pmatrix}
$$

אינה מוגדרת, כי מספר העמודות של המטריצה השמאלית
$(3)$
שונה ממספר השורות של המטריצה הימנית
$(2)$.


\item 
נסתכל על דוגמה שיש לה קשר לתורת הגרפים (שחורגת מהקורס). גרף הוא קבוצת קודקודים שחלקם מחוברים ביניהם ע"י צלעות. לדוגמה:

\begin{figure}[H]
    \centering
    \includegraphics[width=0.4\textwidth]{images/Adjacency.png}
    \caption{ריבוע עם אחד מאלכסוניו}
    \label{fig:Adjacency}
\end{figure}

ניתן לקודד את כל המידע הרלוונטי לגרף זה במטריצה 
$M$
(שנקראת מטריצת השכנויות). אם קיימת צלע בין 
$p_i$
ל-
$p_j$,
נכתוב 
$(M)_{ij}=1$. 
אם לא קיימת צלע ביניהם, נכתוב
$(M)_{ij}=0$. 
כך נקבל
$$
.M =
\begin{pmatrix}
0 & 1 & 1 & 1 \\
1 & 0 & 1 & 0 \\
1 & 1 & 0 & 1 \\
1 & 0 & 1 & 0
\end{pmatrix}
$$

מה המשמעות של כפל? נחשב את 
$M^2=M\cdot M$
ונקבל
$$
.M^2 =
\begin{pmatrix}
3 & 1 & 2 & 1 \\
1 & 2 & 1 & 2 \\
2 & 1 & 3 & 1 \\
1 & 2 & 1 & 2
\end{pmatrix}
$$

אם המטריצה 
$M$
מתארת מספרי מסלולים באורך
$1$, 
אז המטריצה
$M^2$
מתארת מספרי מסלולים באורך
$2$. 
למשל, בין
$p_1$
לעצמו יש
$3$
מסלולים באורך
$2$,
שהם:
$$
p_1\to p_2\to p_1, \quad p_1\to p_3\to p_1,\quad p_1\to p_4\to p_1
$$

זה מתאים לחישוב הבא:

$$
.(M^2)_{11}=\sum_{k=1}^4 (M)_{1k}(M)_{k1}=0\cdot0+1\cdot1+1\cdot1+1\cdot1=3
$$

כל מחובר 
$(M)_{1k}(M)_{k1}$
שתורם 
$1$
לסכום, מתאר מסלול
$p_1\to p_k\to p_1$ 
עם תחנת ביניים
$p_k$.
לכן, הסכום מתאר את מספר המסלולים מ-
$p_1$
לעצמו באורך
$2$.

ובהכללה: 
לכל
$k\in\N$
אפשר לחשב את החזקה
$M^k=\underbrace{M \cdot M \cdot \ldots \cdot M}_{\text{פעמים} \ k}$
ואז
$(M^k)_{ij}$
שווה למספר המסלולים באורך
$k$
מ-
$p_i$
ל-
$p_j$,
כאשר אין הגבלה על תחנות הביניים (יכולה להיות חזרה על קודקוד אחד או יותר).
רעיון זה תקף לכל גרף, והוא דוגמה טובה לשימוש של כפל מטריצות.


\end{enumerate}
\end{example}

\begin{remark}
ראינו כי אין חוק חילוף למטריצות. בהרבה מקרים מתקיים 
$
AB\neq BA
$
גם אם שתי פעולות הכפל מוגדרות. לפעמים כן מתקיים
$AB=BA$,
ואז נאמר שהמטריצות
$A,B$
מתחלפות. ניתקל במקרים כאלה בהמשך, אבל בינתיים חשוב להבין שרוב המטריצות אינן מתחלפות.
\end{remark}

\begin{exercise}
חשבו את
$AB$
עבור
$$
.A=\begin{pmatrix}
1 & -2 & 3 \\
-4 & 5 & -6
\end{pmatrix}, \quad
\begin{pmatrix}
3 & 2 & 1 \\
2 & 3 & 0 \\
3 & -1 & 7
\end{pmatrix}
$$
\end{exercise}

\begin{solution}
נחשב לפי ההגדרה:

$$
AB=
\begin{pmatrix}
1 & -2 & 3 \\
-4 & 5 & -6
\end{pmatrix}
\begin{pmatrix}
3 & 2 & 1 \\
2 & 3 & 0 \\
3 & -1 & 7
\end{pmatrix}
=
\begin{pmatrix}
8 & -7 & 22 \\
-20 & 13 & -46
\end{pmatrix}
$$
\end{solution}

\begin{proposition}\label{prop:product-column}
\addpropositiontolist{כפל מטריצות ככפל המטריצה השמאלית בוקטורי העמודה של המטריצה הימנית}
יהיו
$A\in\M_{m\times k}(\F), B\in\M_{k\times n}(\F)$.
נניח שהעמודות של
$B$
הם הוקטורים
$v_1,v_2,...,v_n\in\F^k$.
אז מתקיים

$$
AB=A\begin{pmatrix}
| & | & & | \\
v_1 & v_2 & \cdots & v_n \\
| & | & & |
\end{pmatrix}=\begin{pmatrix}
| & | & & | \\
Av_1 & Av_2 & \cdots & Av_n \\
| & | & & |
\end{pmatrix}
$$

כלומר התוצאה של כפל ב-
$A$
משמאל הוא שכל וקטור עמודה 
מוכפל ב-
$A$
לחוד.

\end{proposition}

\begin{proof}
ישירות לפי ההגדרות של כפל מטריצוות וכפל מטריצה בוקטור עמודה (שהוא מטריצה עם עמודה אחת). 
לכל
$1\leq j\leq n$
נכתוב
$(B)_{ij}=b_{ij}$, 
כלומר
$$
.v_j=\begin{pmatrix}b_{1j} \\ b_{2j} \\ \vdots \\ b_{kj}\end{pmatrix}
$$

אז לפי הגדרות הכפל נובע כי

$$
.(AB)_{ij}=\sum_{l=1}^k (A)_{il}(B)_{lj}=\sum_{l=1}^k A_{il}b_{lj}=(Av_j)_i
$$

לכן וקטור העמודה ה-
$j$
של
$AB$
הוא אכן
$Av_j$.
\end{proof}

\begin{proposition}\label{prop:matrix-multiplication}
\addpropositiontolist{תכונות של כפל מטריצות}
יהיו
$A_1,A_2\in\M_{m\times n}(\F)$,
$B_1,B_2\in\M_{n\times p}(\F)$,
$C\in\M_{p\times q}(\F)$
ו-
$\alpha\in\F$.
אז מתקיים

\begin{enumerate}[label=\alph*.]
\item 
$$
(A_1+A_2)B_1=A_1B_1+A_2B_1
$$


\item 
$$
A_1(B_1+B_2)=A_1B_1+A_1B_2
$$


\item 
$$
A_1(B_1C)=(A_1B_1)C
$$


\item 
$$
(\alpha A_1)B_1=\alpha(A_1B_1)=A_1(\alpha B_1)
$$


\item 
$$
\0_{m\times n} B_1=\0_{m\times p}=A_1\0_{n\times p}
$$


\end{enumerate}
\end{proposition}

\begin{proof}
רוב הסעיפים נובעים משילוב בין טענה
\ref{prop:matrix-vector}
לטענה
\ref{prop:product-column}.
לכן נסתפק בהוכחת סעיף ג' (חוק הקיבוץ לכפל מטריצות), שהוא היוצא מן הכלל. נניח כי
$$
.(A_1)_{ik} = a_{ik}, \quad
(B_1)_{kl} = b_{kl}, \quad
(C)_{lj} = c_{lj}
$$

אם כן:

$$
(B_1C)_{kj} = \sum_{l=1}^{p} b_{kl} c_{lj},
$$
ולכן
\begin{equation}\label{eq:associativity1}
.(A_1(B_1C))_{ij} = \sum_{k=1}^{n}a_{ik}(B_1C)_{kj}
             = \sum_{k=1}^{n} a_{ik} \left( \sum_{l=1}^{p} b_{kl} c_{lj} \right)
             = \sum_{k=1}^{n} \sum_{l=1}^{p}a_{ik}b_{kl}c_{lj}.
\end{equation}

כעת נחשב את
$((A_1B_1)C)_{ij}$:
$$
(A_1B_1)_{il} = \sum_{k=1}^{n}a_{ik}b_{kl},
$$
ולכן
\begin{equation}\label{eq:associativity2}
.((A_1B_1)C)_{ij} = \sum_{l=1}^{p}(A_1B_1)_{il}c_{l j}
              = \sum_{l=1}^{p} \left( \sum_{k=1}^{n}a_{ik}b_{kl}\right) c_{l j}
              = \sum_{l=1}^{p}\sum_{k=1}^{n} a_{ik}b_{kl}c_{lj}.
\end{equation}

נשים לב כי יש שוויון בין שני הביטויים בצד ימין של
\ref{eq:associativity1}
ו-
\ref{eq:associativity2},
כי לפי חוק החילוף אין חשיבות לסדר הסכימה (אפשר לסכום קודם לפי 
$l$
או קודם לפי 
$k$
בלי לשנות את הסכום). לכן, לכל
$i,j$
מתקיים
$(A_1(B_1C))_{ij}=((A_1B_1)C)_{ij} $, ומכאן
$$
.A_1(B_1C)=(A_1B_1)C
$$
\end{proof}

\begin{definition}
\adddeftolist{מטריצת היחידה $I_n$}
מטריצת היחידה מסדר 
$n\times n$
(אפשר גם לומר מסדר $n$)
מוגדרת להיות
$$
.I_n=I_{n\times n}=
\begin{pmatrix}
1 &0 &\cdots &0\\
0 &1 &\cdots &0\\
\vdots &\vdots &\ddots &\vdots\\
0 &0 &\cdots &1
\end{pmatrix}
$$

זוהי המטריצה שבה כל איבר על האלכסון הראשי הוא
$1$,
וכל שאר האיברים הם
$0$. 
או באופן מתמטי:

$$
\begin{aligned}
\forall 1\leq i\leq n \quad (I_n)_{ii}=1 \\
\forall 1\leq i,j\leq n  \quad i\neq j \implies &(I_n)_{ij}=0
\end{aligned}
$$
\end{definition}

\begin{example}\
\begin{enumerate}[label=\alph*.]
\item 
$$I_1=\begin{pmatrix}1\end{pmatrix}$$

זהו בעצם הסקלר 
$1$,
עם הבדל פורמלי של סוגריים.


\item 
$$
I_2=
\begin{pmatrix}
1 &0 \\
0 &1
\end{pmatrix}
$$


\item 
$$
I_3=
\begin{pmatrix}
1 &0 &0 \\
0 &1 &0 \\
0 &0 &1
\end{pmatrix}
$$


\end{enumerate}
\end{example}

מטריצת היחידה (מכל סדר) משחקת את התפקיד של הסקלר
$1$
במובן של ניטרליות ביחס לכפל. ליתר דיוק:

\begin{proposition}
\addpropositiontolist{תכונות של מטריצת היחידה}
לכל 
$A\in\M_{m\times n}(\F)$
מתקיים
$AI_n=A=I_m A$.
\end{proposition}

\begin{proof}

לכל
$1\leq i,j \leq n$
מתקיים

$$
.(AI_n)_{ij}=\sum_{l=1}^n (A)_{il}(I_n)_{lj}
$$

נשים לב כי לפי הגדרת מטריצת היחידה, מתקיים
$(I)_{lj}=0$
לכל
$l\neq j$.
אז רק איבר אחד בסכום באגף ימין לא מתאפס, והוא מתאים להצבה
$l=j$. 
לכן

$$
.(AI_n)_{ij}=0+...+0+(A)_{ij}(I_n)_{jj}+0+...+0=(A)_{ij}
$$

אז 
$AI_n=A$
כי יש שוויון בין כל האיברים המתאימים. באופן דומה:

$$
.(I_m A)_{ij}=\sum_{l=1}^m (I_m)_{il}A_{lj}=0+...+0+(I)_{ii}(A)_{ij}+0+...+0=(A)_{ij}
$$

ולכן
$I_m A=A$.

\end{proof}

\begin{remark}
אם
$m=n$,
אז לפי הטענה מתקיים
$AI_n=A=I_nA$.
במילים: 
$I_n$
מתחלפת עם כל מטריצה מסדר 
$n\times n$.
\end{remark}

\subsection{שחלוף מטריצות}
ניתן לעשות המרה בין שורות לעמודות. פעולה זו נקראת שחלוף.

\begin{definition}
\adddeftolist{שחלוף}
תהי
$A\in\M_{m\times n}(\F)$. 
נגדיר את המטריצה המשוחלפת
$A^t\in\M_{n\times m}(\F)$
ע"י
$$
(A^t)_{ij}=A_{ji}
$$
לכל
$1\leq i\leq n,1\leq j\leq m$.
\end{definition}

\begin{remark}
$A^t$
זה סימון מיוחד, לא סימון חזקה. כאן
$t$
זה מלשון
transpose
(שחלוף באנגלית).
\end{remark}

\begin{example}\
\begin{enumerate}[label=\alph*.]
\item 
עבור וקטור שורה
$v=(\alpha_1,\alpha_2,...,\alpha_n)$
נקבל
$$
.v^t=\begin{pmatrix}
\alpha_1\\
\alpha_2\\
\vdots \\
\alpha_n
\end{pmatrix}
$$

\item 
עבור
$$
A=\begin{pmatrix}
a &b\\
c & d
\end{pmatrix}
$$
נקבל
$$
.A^t=\begin{pmatrix}
a &c\\
b & d
\end{pmatrix}
$$


\item 
עבור
$$
B=\begin{pmatrix}
1 &2\\
3 & 4\\
5 & 6
\end{pmatrix}
$$
נקבל
$$
.B^t=\begin{pmatrix}
1 &3 &5\\
2 & 4 &6
\end{pmatrix}
$$


\end{enumerate}
\end{example}


\begin{proposition}\label{prop:transpose}
\addpropositiontolist{תכונות של שחלוף}
לכל 
$A,B\in\M_{m\times n}(\F), C\in\M_{n\times k}(\F)$
ולכל
$\alpha\in\F$
מתקיים:
\begin{enumerate}[label=\alph*.]
\item 
$$
(A^t)^t=A
$$


\item 
$$
(A+B)^t=A^t+B^t
$$

\item 
$$
(\alpha A)^t=\alpha A^t
$$


\item 
$$
(AC)^t=C^tA^t
$$


\end{enumerate}
\end{proposition}

\begin{proof}
נראה שוויון בין האיברים הכלליים של המטריצות בשני האגפים של כל סעיף.
\begin{enumerate}[label=\alph*.]
\item 
שתי המטריצות מסדר
$m\times n$
ולכל 
$1\leq i\leq m, 1\leq j\leq n$
מתקיים
$$
.((A^t)^t)_{ij}=(A^t)_{ji}=(A)_{ij}
$$


\item 
שתי המטריצות מסדר
$n\times m$
ולכל 
$1\leq i\leq n, 1\leq j\leq m$
מתקיים
$$
.(A+B)^t_{ij}=(A+B)_{ji}=(A)_{ji}+(B)_{ji}=(A^t)_{ij}+(B^t)_{ij}=(A^t+B^t)_{ij}
$$


\item
שתי המטריצות מסדר
$n\times m$
ולכל 
$1\leq i\leq n, 1\leq j\leq m$
מתקיים
$$
.((\alpha A)^t)_{ij}=(\alpha A)_{ji}=\alpha (A)_{ji}=\alpha (A^t)_{ij}=(\alpha A^t)_{ij}
$$


\item 
כאן
$AC$
מסדר 
$m\times k$
ולכן
$(AC)^t$
מסדר
$k\times m$.
באופן דומה,
$C^t$
מסדר 
$k\times n$
ו-
$A^t$
מסדר
$n\times m$. 
אז גם
$C^tA^t$
מסדר
$k\times m$,
ולכל
$1\leq i\leq k, 1\leq j\leq m$
מתקיים:
$$
\begin{aligned}
&((AC)^t)_{ij}=(AC)_{ji}=\sum_{l=1}^n (A)_{jl}(C)_{li}=\sum_{l=1}^n (A^t)_{lj}(C^t)_{il}=\sum_{l=1}^n (C^t)_{il}(A^t)_{lj} \\
&=(C^tA^t)_{ij}
\end{aligned}
$$
לקראת הסוף השתמשנו בחוק החילוף כדי להתאים את סדר הכפל בתוך הסכום לתבנית של 
$C^tA^t$.


\end{enumerate}
\end{proof}

\begin{remark}
בתנאי הטענה, המטריצה
$A^tC^t$
מוגדרת אם ורק אם מספר העמודות של
$A^t$
שווה למספר השורות של
$C^t$,
כלומר אם ורק אם 
$m=k$.
במקרה זה הסדר של
$A^tC^t$
הוא
$n\times n$. 
לעומת זאת, הסדר של
$(AC)^t$
הוא
$k\times m$,
או
$m\times m$
תחת ההנחה ש-
$A^tC^t$
אכן מוגדרת. גם במקרה
$m=k=n$
אין הכרח שיתקיים
$A^tC^t=C^tA^t$.
ולסיכום: אין הכרח שיתקיים
$(AC)^t=A^tC^t$.
\end{remark}

\begin{example}
ניקח
$$
A = \begin{pmatrix}
1 & 2 & 3 \\
4 & 5 & 6
\end{pmatrix}, \quad
C = \begin{pmatrix}
1 & 0 \\
-1 & 2 \\
0 & 3
\end{pmatrix}.
$$

נחשב גם את
$(AC)^t$
וגם את
$C^tA^t$
ונראה שאכן יש שוויון (כצפוי). ראשית, נחשב את
$AC$:

$$
AC=
\begin{pmatrix}
1 & 2 & 3 \\
4 & 5 & 6
\end{pmatrix}
\begin{pmatrix}
1 & 0 \\
-1 & 2 \\
0 & 3
\end{pmatrix}
=
\begin{pmatrix}
-1 & 13 \\
-1 & 28
\end{pmatrix}
$$

לכן

$$
.(AC)^t=\begin{pmatrix}
-1 & -1 \\
13 & 28
\end{pmatrix}
$$

נעבור לחישוב
$C^tA^t$:

$$
C^t A^t=
\begin{pmatrix}
1 & -1 & 0 \\
0 & 2 & 3
\end{pmatrix}
\begin{pmatrix}
1 & 4 \\
2 & 5 \\
3 & 6
\end{pmatrix}
=
\begin{pmatrix}
-1 & -1 \\
13 & 28
\end{pmatrix}
$$

אז אכן קיבלנו 
$(CA)^t=C^tA^t$.
\end{example}

\begin{exercise}
עבור
$$
A = \begin{pmatrix}
1+i & 2 \\
3 & 4-i
\end{pmatrix}, \quad
B = \begin{pmatrix}
2 & i \\
1-i & 3
\end{pmatrix}
$$

הראו ע"י חישוב ישיר כי
$(AB)^t=B^tA^t\neq A^tB^t$.
\end{exercise}

\begin{solution}
נחשב את
$AB$ 
ונקבל

$$
.AB=
\begin{pmatrix}
4 & 5 + i \\
9 - 5i & 12
\end{pmatrix}
$$

לכן
$$
.(AB)^t=\begin{pmatrix}
4 & 9-5i \\
5+i & 12
\end{pmatrix}
$$

נעבור לחישוב השני:

$$
\begin{aligned}
B^t A^t &=
\begin{pmatrix}
2 & 1-i \\
i & 3
\end{pmatrix}
\begin{pmatrix}
1+i & 3 \\
2 & 4-i
\end{pmatrix}
=
\begin{pmatrix}
4 & 9 - 5i \\
5 + i & 12
\end{pmatrix}
\end{aligned}
$$

אז אכן יש שוויון. לעומת זאת:

$$
A^tB^t =
\begin{pmatrix}
1+i & 3 \\
2 & 4-i
\end{pmatrix}
\begin{pmatrix}
2 & 1-i \\
i & 3
\end{pmatrix}
=
\begin{pmatrix}
2+5i & 11 \\
5 + 4i & 14-5i
\end{pmatrix}\neq B^tA^t
$$
\end{solution}

\section{מטריצות מיוחדות}

נרצה להגדיר קבוצות של מטריצות מיוחדות, שיופיעו בהמשך הקורס.

\subsection{מטריצות ריבועיות}

מטריצות ריבועיות הן סוג נפוץ של מטריצות שניתקל בהן לאורך הקורס. כשמן כן הן - מטריצות שנראות כמו ריבוע עם מספרים בתוכו. הרבה מטריצות בקורס יהיו ריבועיות.

\begin{definition}
\adddeftolist{מטריצה ריבועית}
מטריצה 
$A$
נקראת ריבועית אם מספר השורות שלה שווה למספר העמודות, כלומר אם קיים 
$n\in\N$
עבורו
$A\in\M_{n\times n}(\F)$.
\end{definition}

\begin{example}
המטריצות הבאות הן ריבועיות:

$$
\begin{pmatrix}
1 &2 \\
3 &4
\end{pmatrix}, \quad
\begin{pmatrix}
i &2 &3 \\
4 &0 &6 \\
7 &8 &i
\end{pmatrix}
$$

המטריצות הבאות אינן ריבועיות:

$$
\begin{pmatrix}
1 &2 &3\\
4 &5 &6
\end{pmatrix}, \quad
\begin{pmatrix}
1 &2  \\
-2 &-1  \\
0 &5 \\
1 & 10
\end{pmatrix}, \quad
\begin{pmatrix}
i \\
1  \\
0 
\end{pmatrix}
$$
\end{example}

\begin{remark}
אם 
$A,B\in\M_{n\times n}(\F)$
אז גם
$A\pm B,AB,BA\in\M_{n\times n}(\F)$. 
במילים: הקבוצה 
$\M_{n\times n}(F)$
סגורה לחיבור, חיסור וכפל.
\end{remark}

\begin{definition}\label{def:matrix-polynomial}
\adddeftolist{חזקות ופולינום של מטריצה}
תהי 
$A\in\M_{n\times n}(\F)$.
אז נגדיר 

$$
A^0=I_n, \quad A^1=A, \quad A^2=A\cdot A, \quad A^3=A\cdot A\cdot A
$$

ולכל
$k\in\N$

$$
.A^k = \underbrace{A \cdot A \cdot \ldots \cdot A}_{\text{פעמים} \ k}
$$

עבור פולינום
$$
p(x)=a_mx^m+a_{m-1}x^{m-1}+...+a_1x+a_0
$$
נגדיר את ההצבה של
$A$
בו ע"י
$$
.p(A)=a_mA^m+a_{m-1}A^{m-1}+...+a_1A+a_0I_n
$$
\end{definition}

\begin{example}
נחשב את
$p(A)$
עבור
$$
A=\begin{pmatrix}
1 &2 \\
0 &-1
\end{pmatrix}
$$
והפולינום
$$
.p(x)=x^3+2x^2+x+1
$$

לצורך כך צריך לחשב את החזקות
$A^3, A^2$:

$$
\begin{aligned}
A^2&=\begin{pmatrix}
1 &2 \\
0 &-1
\end{pmatrix}\begin{pmatrix}
1 &2 \\
0 &-1
\end{pmatrix}=\begin{pmatrix}
1 &0 \\
0 &1
\end{pmatrix}=I_2\\
\implies A^3&=A^2\cdot A=I_2 A=A \\
\implies p(A)&=A^3+2A^2+A+I_2=2A+3I_2=2\begin{pmatrix}
1 &2 \\
0 &-1
\end{pmatrix}+3\begin{pmatrix}
1 &0 \\
0 &1
\end{pmatrix} \\
&=\begin{pmatrix}
5 &4 \\
0 &1
\end{pmatrix}
\end{aligned}
$$

\end{example}

\begin{remark}
הדוגמה האחרונה גם מראה שייתכן
$A^2=I_n$
גם אם
$A\neq\pm I_n$.
אמנם לסקלר יש שני שורשים (או אחד במקרה של
$0$),
אבל למטריצת יחידה מסדר לפחות 
$2$
יש אינסוף שורשים. בקורס שלנו אסור להוציא שורש של מטריצה כי זו לא פעולה שנגדיר (לפעמים יש אינסוף שורשים, לפעמים יש מספר סופי כלשהו של שורשים, ולפעמים אין שורש).

באופן דומה, מתקיים
$B^2=\0_{2\times 2}$
עבור
$$
.B=\begin{pmatrix}
0 &1 \\
0 &0
\end{pmatrix}
$$
אז לא רק שקיימת למטריצת האפס
$\0_{2\times 2}$
שורש שאינו עצמה, זה גם מראה שכלל הצמצום למטריצות לא מתקיים באופן כללי.

$$
B\cdot B=B\cdot \0_{n\times n} \nRightarrow B=\0_{n\times n}
$$

לא ניתן "לחלק" במטריצה, אלא אם כן היא מסוג מסוים שעוד לא הגדרנו (מטריצה הפיכה). זה לא המקרה פה.

\end{remark}

\subsection{מטריצות סימטריות ומטריצות אנטי-סימטריות}

כאן הכוונה לסימטריה בין שורות לעמודות. לכן יש קשר לפעולת השחלוף.

\begin{definition}
\adddeftolist{מטריצה סימטרית ומטריצה אנטי-סימטרית}
תהי 
$A\in\M_{n\times n}(\F)$
מטריצה ריבועית.
\begin{enumerate}[label=\alph*.]
\item
$A$
נקראת סימטרית אם מתקיים
$A^t=A$,
או באופן שקול:
$$
\forall 1\leq i,j\leq n \quad (A)_{ji}=(A)_{ij}
$$


\item 
$A$
נקראת אנטי-סימטרית אם מתקיים
$A^t=-A$,
או באופן שקול:
$$
\forall 1\leq i,j\leq n \quad (A)_{ji}=-(A)_{ij}
$$

\end{enumerate}
\end{definition}

\begin{example}\
\begin{enumerate}[label=\alph*.]
\item
המטריצות הבאות הן סימטריות:
$$
\begin{pmatrix}
1 &2\\
2 &3
\end{pmatrix}, \begin{pmatrix}
1 &i &0\\
i &3 &-1  \\
0 &-1 &5
\end{pmatrix}
$$


\item
המטריצות הבאות הן אנטי-סימטריות:
$$
\begin{pmatrix}
0 &1\\
-1 &0
\end{pmatrix}, \begin{pmatrix}
0 &i &-2\\
-i &0 &-1  \\
2 &1 &0
\end{pmatrix}
$$


\end{enumerate}
\end{example}

\begin{proposition}\label{prop:sym-antisym}
\addpropositiontolist{תכונות של מטריצות אנטי-סימטריות ומטריצות סימטריות}
תהי
$P\in\M_{n\times n}(\F)$
מטריצה ריבועית.
\begin{enumerate}[label=\alph*.]
\item
אם
$P$
אנטי-סימטרית, אז כל איבריה על האלכסון הראשי שווים ל-
$0$.


\item 
אם
$P$
סימטרית וגם אנטי-סימטרית, אז
$P=\0_{n\times n}$.


\item 
לכל
$B\in\M_{n\times n}(\F)$
קיימות
$S\in\M_{n\times n}(\F)$
סימטרית ו-
$A\in\M_{n\times n}(\F)$
אנטי-סימטרית כך שמתקיים
$B=S+A$.


\end{enumerate}
\end{proposition}

\begin{proof}\
\begin{enumerate}[label=\alph*.]
\item 
$$
\forall 1\leq i\leq n\quad(P)_{ii}=-(P)_{ii} \implies (P)_{ii}=0
$$
כלומר יש רק אפסים על האלכסון הראשי של
$P$.


\item
מצד אחד
$
P^t=P
$
ומצד שני
$
.P^t=-P
$
לכן

$$
P=-P \implies P+P=-P+P \implies 2P=\0_{n\times n} \implies P=\0_{n\times n}
$$

למעשה, השתמשנו באלגברה רגילה כדי לקבל 
$P=\0_{n\times n}$
כמו עבור משוואה במשתנה סקלרי. מותר לחלק בסקלר
$2$,
שזה כמו להכפיל בסקלר
$\frac{1}{2}$.


\item 
ניקח
$$
.S=\frac{1}{2}(B+B^t), \quad A=\frac{1}{2}(B-B^t)
$$
אז לפי טענה
\ref{prop:transpose}
מתקיים
$$
\begin{aligned}
S^t&=\frac{1}{2}(B+B^t)^t=\frac{1}{2}(B^t+B)=S \\
A^t&=\frac{1}{2}(B-B^t)^t=\frac{1}{2}(B^t-B)=-A \\
B&=\frac{1}{2}(B+B^t)+\frac{1}{2}(B-B^t)=S+A
\end{aligned}
$$

לכן
$S$
סימטרית ו-
$A$
אנטי-סימטרית כנדרש.


\end{enumerate}
\end{proof}

\begin{example}
נראה כי מכפלה של מטריצות סימטריות 
$A,B$
אינה בהכרח סימטרית. אם היינו מנסים להוכיח שהמכפלה היא סימטרית, היינו נתקעים:

$$
(AB)^t=B^tA^t=BA
$$

הבעיה היא שלא בהכרח מתקיים 
$BA=AB$.
אז בתור דוגמה נגדית נמצא שתי מטריצות סימטריות מסדר
$2\times 2$
שאינן מתחלפות. למשל:

$$
A=\begin{pmatrix}
1 &2 \\
2 &0
\end{pmatrix}, \quad
B=\begin{pmatrix}
1 &1 \\
1 &0
\end{pmatrix}
$$

נקבל

$$
AB=\begin{pmatrix}
1 &2 \\
2 &0
\end{pmatrix}\begin{pmatrix}
1 &1 \\
1 &0
\end{pmatrix}=\begin{pmatrix}
3 &1 \\
2 &2
\end{pmatrix}
$$

וזו לא מטריצה סימטרית.

\end{example}

\begin{exercise}
הסבירו מדוע כל מטריצה אנטי-סימטרית
$A\in\M_{2\times 2}(\F)$
היא מהצורה 
$$
\begin{pmatrix}
0 &a\\
-a &0
\end{pmatrix}
$$
עבור 
$a\in\F$.
הסיקו שלכל שתי מטריצות אנטי-סימטריות 
$A,B\in\M_{2\times 2}(\F)$
שאינן 
$\0_{2\times 2}$
מתקיים
$AB=BA$
ומכפלה זו אינה אנטי-סימטרית.
\end{exercise}

\begin{solution}
האיברים על האלכסון הראשי של כל מטריצה אנטי-סימטרית הם
$0$.
נסמן
$(A)_{12}=a$
ונקבל
$(A)_{21}=-a$.
לכן 
$$
,A=\begin{pmatrix}
0 &a\\
-a &0
\end{pmatrix}
$$
ובאופן דומה קיים
$b\in\F$
כך שמתקיים
$$
.B=\begin{pmatrix}
0 &b\\
-b &0
\end{pmatrix}
$$

נחשב מכפלות:
$$
.AB=\begin{pmatrix}
0 &a\\
-a &0
\end{pmatrix}\begin{pmatrix}
0 &b\\
-b &0
\end{pmatrix}=\begin{pmatrix}
-ab &0 \\
0 &-ab
\end{pmatrix}=\begin{pmatrix}
0 &b\\
-b &0
\end{pmatrix}\begin{pmatrix}
0 &a\\
-a &0
\end{pmatrix}=BA
$$

זו לא מטריצה אנטי-סימטרית כי 
$ab\neq 0$
לפי ההנחה
$A,B\neq\0$.

\end{solution}

\begin{exercise}\
\begin{enumerate}[label=\alph*.]
\item
לכל 
$A,B\in\M_{n\times n}(\F)$
סימטריות ולכל
$\alpha\in\F$
מתקיים שהמטריצות
$A+B,\alpha A$
הן סימטריות.


\item 
לכל 
$A,B\in\M_{n\times n}(\F)$
אנטי-סימטריות ולכל
$\alpha\in\F$
מתקיים שהמטריצות
$A+B,\alpha A$
הן אנטי-סימטריות.


\end{enumerate}
\end{exercise}

\begin{solution}
\begin{enumerate}[label=\alph*.]
\item
לכל
$A,B\in\M_{n\times n}(\F)$
סימטריות ולכל
$\alpha\in\F$
נקבל לפי טענה
\ref{prop:transpose}
$$
\begin{aligned}
(A+B)^t&=A^t+B^t=A+B \\
(\alpha A)^t&=\alpha A^t=\alpha A
\end{aligned}
$$
ולכן אלו מטריצות סימטריות.

\item 
הפעם נקבל

$$
\begin{aligned}
(A+B)^t&=A^t+B^t=-A-B=-(A+B) \\
(\alpha A)^t&=\alpha A^t=\alpha(-A)=-(\alpha A)
\end{aligned}
$$
ולכן אלו מטריצות אנטי-סימטריות.


\end{enumerate}
\end{solution}

\subsection{מטריצות אלכסוניות ומטריצות משולשיות}

\begin{definition}
\adddeftolist{מטריצה אלכסונית ומטריצה משולשית עליונה/תחתונה}
מטריצה 
$A\in\M_{n\times n}(\F)$
נקראת:

\begin{enumerate}[label=\alph*.]
\item 
אלכסונית אם לכל
$i\neq j$
מתקיים
$(A)_{ij}=0$,
כלומר כל האיברים מחוץ לאלכסון הראשי שווים ל-
$0$.


\item
משולשית עליונה אם לכל
$i>j$
מתקיים
$(A)_{ij}=0$,
כלומר כל האיברים מתחת לאלכסון הראשי שווים ל-
$0$.


\item
משולשית תחתונה אם לכל
$i<j$
מתקיים
$(A)_{ij}=0$,
כלומר כל האיברים מעל לאלכסון הראשי שווים ל-
$0$.


\end{enumerate}
\end{definition}

\begin{example}\
\begin{enumerate}[label=\alph*.]
\item 
המטריצות הבאות הן אלכסוניות:
$$
\begin{pmatrix}
1 &0 \\
0 &2
\end{pmatrix}, \begin{pmatrix}
-1 &0 &0 \\
0 &2 &0 \\
0 &0 &i
\end{pmatrix}, \begin{pmatrix}
4 &0 &0 &0 \\
0 &6 &0 &0 \\
0 &0 &8 &0 \\
0 &0 &0 &0
\end{pmatrix}
$$


\item
המטריצות הבאות הן משולשיות עליונות:
$$
\begin{pmatrix}
1 &3 \\
0 &2
\end{pmatrix}, \begin{pmatrix}
-1 &-i &3 \\
0 &2 &7 \\
0 &0 &i
\end{pmatrix}, \begin{pmatrix}
4 &5 &20 &0 \\
0 &6 &7 &13 \\
0 &0 &8 &5 \\
0 &0 &0 &0
\end{pmatrix}
$$


\item
המטריצות הבאות הן משולשיות תחתונות:
$$
\begin{pmatrix}
1 &0 \\
-4 &2
\end{pmatrix}, \begin{pmatrix}
-1 &0 &0 \\
1 &2 &0 \\
3 &2i &i
\end{pmatrix}, \begin{pmatrix}
4 &0 &0 &0 \\
-2 &6 &0 &0 \\
3 &7 &8 &0 \\
4 &-5 &6 &0
\end{pmatrix}
$$


\end{enumerate}
\end{example} 

\begin{remark}
נשים לב שלפי ההגדרות, לכל 
$A$
ריבועית מתקיים:
\begin{enumerate}[label=\alph*.]
\item 
$A$
משולשית עליונה וגם משולשית תחתונה 
$\iff$
$A$
אלכסונית.


\item 
$A$
משולשית עליונה
$\iff$
$A^t$
משולשית תחתונה.


\item
$A$
משולשית תחתונה
$\iff$
$A^t$
משולשית עליונה.


\end{enumerate}
\end{remark}

\begin{proposition}
\addpropositiontolist{סגירות לכפל של מטריצות משולשיות ומטריצות אלכסוניות}
יהיו
$A,B\in\M_{n\times n}(\F)$.
\begin{enumerate}[label=\alph*.]
\item
אם
$A,B$
משולשיות עליונות, אז
$AB$
גם משולשית עליונה.


\item
אם
$A,B$
משולשיות תחתונות, אז
$AB$
גם משולשית תחתונה.


\item
אם
$A,B$
אלכסוניות, אז
$AB$
גם אלכסונית.


\end{enumerate}
\end{proposition}

\begin{proof}\
\begin{enumerate}[label=\alph*.]
\item 
לכל 
$j<i$
נסתכל על האיבר הכללי של המכפלה:

$$
(AB)_{ij}=\sum_{l=1}^n (A)_{il}(B)_{lj}
$$

עבור 
$l< i$
מתקיים
$(A)_{il}=0$
כי 
$A$
משולשית עליונה. אז 
$i-1$
המכפלות הראשונות בסכום מתאפסות.
עבור 
$l\geq i>j$
מתקיים
$(B)_{lj}=0$
כי 
$B$
משולשית עליונה. אז גם המכפלות הנותרות בסכום מתאפסות, ולכן
$(AB)_{ij}=0$
כנדרש.


\item 
אפשר לחקות את ההוכחה של הסעיף הקודם, אבל יש קיצור דרך: אם 
$A,B$
משולשיות תחתונות, אז
$A^t,B^t$
משולשיות עליונות. לפי הסעיף הקודם, נובע כי
$B^tA^t=(AB)^t$
משולשית עליונה ולכן
$AB$
משולשית תחתונה.


\item 
שוב אפשר להשתמש בסעיפים הקודמים. אם 
$A,B$
אלכסוניות, אז הן גם משולשיות עליונות וגם משולשיות תחתונות. לכן, לפי הסעיפים הקודמים
$AB$
גם משולשית עליונה וגם משולשית תחתונה, כלומר היא אלכסונית.


\end{enumerate}
\end{proof}


\begin{remark}
לאחר שהבנו שהמכפלה של שתי מטריצות אלכסוניות (מאותו הסדר) היא בהכרח אלכסונית, קל להשתכנע שהנוסחה הבאה נכונה:

נסמן
$$
.A=\begin{pmatrix}
a_1 &0 &0 &\cdots &0 \\
0 &a_2 &0 &\cdots &0 \\
0 &0 &a_3 &\cdots &0 \\
\vdots &\vdots &\vdots &\ddots &\vdots \\
0 &0 &0 &\cdots &a_n
\end{pmatrix}, \quad
B=\begin{pmatrix}
b_1 &0 &0 &\cdots &0 \\
0 &b_2 &0 &\cdots &0 \\
0 &0 &b_3 &\cdots &0 \\
\vdots &\vdots &\vdots &\ddots &\vdots \\
0 &0 &0 &\cdots &b_n
\end{pmatrix}
$$
אז לפי חישוב איברי האלכסון הראשי של המכפלה, נקבל
$$
.AB=\begin{pmatrix}
a_1b_1 &0 &0 &\cdots &0 \\
0 &a_2b_2 &0 &\cdots &0 \\
0 &0 &a_3b_3 &\cdots &0 \\
\vdots &\vdots &\vdots &\ddots &\vdots \\
0 &0 &0 &\cdots &a_nb_n
\end{pmatrix}
$$

בפרט, אם
$B=A$
נקבל
$$
.A^2=\begin{pmatrix}
a_1^2 &0 &0 &\cdots &0 \\
0 &a_2^2 &0 &\cdots &0 \\
0 &0 &a_3^2 &\cdots &0 \\
\vdots &\vdots &\vdots &\ddots &\vdots \\
0 &0 &0 &\cdots &a_n^2
\end{pmatrix}
$$
נכפיל שוב ב-
$A$
ונקבל:
$$
.A^3=\begin{pmatrix}
a_1^3 &0 &0 &\cdots &0 \\
0 &a_2^3 &0 &\cdots &0 \\
0 &0 &a_3^3 &\cdots &0 \\
\vdots &\vdots &\vdots &\ddots &\vdots \\
0 &0 &0 &\cdots &a_n^3
\end{pmatrix}
$$
ובהכללה (או באינדוקציה למי שמכיר) לכל
$k\in\N$
נובע כי
$$
.A^k=\begin{pmatrix}
a_1^k &0 &0 &\cdots &0 \\
0 &a_2^k &0 &\cdots &0 \\
0 &0 &a_3^k &\cdots &0 \\
\vdots &\vdots &\vdots &\ddots &\vdots \\
0 &0 &0 &\cdots &a_n^k
\end{pmatrix}
$$

אז קל יחסית לחשב חזקות של מטריצה אלכסונית, בניגוד למטריצה כללית עבורה יש הרבה יותר חישובים (בוודאי אם 
$n$
והמעריך
$k$
הם מספרים גדולים). 

\end{remark}

\begin{definition}
\adddeftolist{מטריצה סקלרית}
מטריצה 
$A\in\M_{n\times n}(\F)$
נקראת סקלרית אם קיים
$\alpha\in\F$
כך ש-
$A=\alpha I_n$.
\end{definition}

\begin{exercise}
הראו כי לכל מטריצה סקלרית
$A=\alpha I_n$
ולכל מטריצה
$B\in\M_{n\times n}(\F)$
מתקיים
$AB=BA=\alpha B$.
\end{exercise}

\begin{solution}
יהיו 
$A,B$
כנ"ל. אז מתקיים

$$
,AB=(\alpha I_n)B=\alpha(I_nB)=\alpha B
$$

ובנוסף

$$
.BA=B(\alpha I_n)=\alpha(BI_n)=\alpha B
$$

לכן יש שוויון.

\end{solution}

\subsection{מטריצות אלמנטריות}
כאן יש קשר הדוק לפעולות שורה אלמנטריות. נשתמש בראשי התיבות פש"א עבור "פעולת שורה אלמנטרית".

\begin{definition}
\adddeftolist{מטריצה אלמנטרית}
מטריצה
$A\in\M_{n\times n}(\F)$
נקראת אלמנטרית אם היא מתקבלת ממטריצת היחידה
$I_n$
ע"י פש"א אחת.

עבור פש"א 
$S$,
נסמן ב-
$S(I_n)$
את המטריצה האלמנטרית המתאימה שמתקבלת מ-
$I_n$
ע"י 
$S$.
באופן כללי, לכל מטריצה
$A$
נסמן ב-
$S(A)$
את המטריצה שמתקבלת מ-
$A$
ע"י
$S$.
\end{definition}

\begin{example}\
\begin{enumerate}[label=\alph*.]
\item 
המטריצות הבאות הן אלמנטריות כי הן מתאימות להחלפה בין שורות של מטריצת יחידה:

$$
\begin{aligned}
S=(R_1\leftrightarrow R_2) \implies S(I_2)&=\begin{pmatrix}
0 &1 \\
1 & 0
\end{pmatrix} \\
S=(R_1\leftrightarrow R_3)\implies S(I_3)&=\begin{pmatrix}
0 &0 &1 \\
0 & 1 &0 \\
1 &0 &0
\end{pmatrix} \\
S=(R_2\leftrightarrow R_4) \implies S(I_4)&=\begin{pmatrix}
1 &0 &0 &0 \\
0 & 0 &0 &1 \\
0 &0 &1 &0 \\
0 &1 &0 &0
\end{pmatrix}
\end{aligned}
$$


\item 
המטריצות הבאות הן אלמנטריות כי הן מתאימות לכפל בסקלר של שורה אחת של מטריצת יחידה:

$$
\begin{aligned}
S=(2R_1\to R_1)\implies S(I_2)&=\begin{pmatrix}
2 & 0\\
0 & 1
\end{pmatrix} \\
S=(3R_2\to R_2)\implies S(I_3)&=\begin{pmatrix}
1 &0 &0 \\
0 & 3 &0 \\
0 &0 &1
\end{pmatrix} \\
S=(4R_3\to R_3)\implies S(I_4)&=\begin{pmatrix}
1 &0 &0 &0 \\
0 & 1 &0 &0 \\
0 &0 &4 &0 \\
0 &0 &0 &1
\end{pmatrix}
\end{aligned}
$$


\item 
המטריצות הבאות הן אלמנטריות כי הן מתאימות להוספת כפולה בסקלר של שורה אחת לשורה אחרת במטריצת יחידה:

$$
\begin{aligned}
S=(R_2+2R_1\to R_2)\implies S(I_2)&=\begin{pmatrix}
1 & 0\\
2 & 1
\end{pmatrix} \\
S=(R_2+3R_1\to R_2)\implies S(I_3)&=\begin{pmatrix}
1 &0 &0 \\
3 & 1 &0 \\
0 &0 &1
\end{pmatrix} \\
S=(R_3+4R_2\to R_3)\implies S(I_4)&=\begin{pmatrix}
1 &0 &0 &0 \\
0 & 1 &0 &0 \\
0 &4 &1 &0 \\
0 &0 &0 &1
\end{pmatrix}
\end{aligned}
$$

\item 
המטריצות הבאות
\underline{אינן}
אלמנטריות כי הן מתקבלות ממטריצת יחידה ע"י יותר מפש"א אחת:

$$
\begin{pmatrix}
2 & 0\\
0 & 3
\end{pmatrix}, \begin{pmatrix}
1 &0 &0 \\
1 & 1 &0 \\
1 &0 &1
\end{pmatrix}, \begin{pmatrix}
0 &0 &0 &1 \\
0 & 0 &1 &0 \\
0 &1 &0 &0 \\
1 &0 &0 &0
\end{pmatrix}
$$
למעשה, במטריצה אלמנטרית רק שורה אחת שונה מהשורה המתאימה ב-
$I_n$,
אלא אם כן זו מטריצה שמתאימה להחלפת שורות. אבל גם אז יש הבדל רק בשתי שורות וקל לזהות את ההחלפה, בניגוד למטריצה השמאלית לעיל שהיא בפירוש לא קשורה להחלפה.

\end{enumerate}
\end{example}

\begin{proposition}\label{prop:elementary}
\addpropositiontolist{פש"א ככפל משמאל במטריצה אלמנטרית}
יהיו
$A\in\M_{m\times n}(\F)$
ו-
$S$
פש"א. אז מתקיים
$S(I_m)A=S(A)$.
\end{proposition}

לפני החישובים הכלליים הדרושים להוכחה, נסתכל על דוגמאות:

\begin{example}\
\begin{enumerate}[label=\alph*.]
\item
$$
\begin{pmatrix}
0 &0 &1 \\
0 & 1 &0 \\
1 &0 &0
\end{pmatrix}\begin{pmatrix}
a &b &c \\
d & e &f \\
g &h &r
\end{pmatrix}=\begin{pmatrix}
g &h &r \\
d & e &f \\
a &b &c 
\end{pmatrix}
$$


\item 
$$
\begin{pmatrix}
1 &0 &0 \\
0 & 3 &0 \\
0 &0 &1
\end{pmatrix}\begin{pmatrix}
a &b &c \\
d & e &f \\
g &h &r
\end{pmatrix}=\begin{pmatrix}
a &b &c \\
3d & 3e &3f \\
g &h &r
\end{pmatrix}
$$


\item 

$$
\begin{pmatrix}
1 &0 &0 \\
3 & 1 &0 \\
0 &0 &1
\end{pmatrix}\begin{pmatrix}
a &b &c \\
d & e &f \\
g &h &r
\end{pmatrix}=\begin{pmatrix}
a &b &c \\
3a+d & 3b+e &3c+f \\
g &h &r
\end{pmatrix}
$$


\end{enumerate}
\end{example}

\begin{proof}
נסמן
$E=S(I_m)$
עם איברים
$e_{ij}=(E)_{ij}$.
נראה ש-
$EA$
מתקבלת מ-
$A$
בדיוק ע"י הפעלת $S$
על שורות
$A$. 
נחלק לשלושה מקרים:

\begin{enumerate}[label=\alph*.]
\item
$S=(R_i\leftrightarrow R_j)$

$E$
מזכירה את מטריצת היחידה, אבל יש ארבעה הבדלים:
$$
e_{ii}=e_{jj}=0, \quad e_{ij}=e_{ji}=1
$$

נחשב:
$$
(EA)_{rk} = \sum_{l=1}^m e_{rl} a_{lk}.
$$
אם
$r \notin \{i,j\}$,
אז $e_{rk}$
הוא בדיוק כמו במטריצת היחידה:
$1$ 
עבור 
$r=k$,
אחרת
$0$.
לכן השורה ה-
$r$
נשארת כפי שהיא.  
לעומת זאת, אם
$r=i$
נקבל
$(EA)_{ik} = a_{jk}$.
אם
$r=j$
נקבל
$(EA)_{jk} = a_{ik}$.  
כלומר החלפנו בין השורות ה-
$i,j$
של 
$A$,
ולכן קיבלנו את
$S(A)$.


\item
$S=(\alpha R_i\to R_i)$
עבור
$\alpha\neq 0$.

כאן
$E$
היא כמעט כמו מטריצת היחידה, עם הבדל אחד:
$
.e_{ii} = \alpha
$

נשתמש בנוסחה מפוצלת כי יש פה שני מקרים:

$$
(EA)_{rk} = \sum_{l=1}^m e_{rl} a_{lk} =
\begin{cases}
a_{rk}, & r \neq i\\
\alpha a_{ik}, & r = i
\end{cases}
$$
כלומר השורה ה-
$i$
של
$A$
הוכפלה ב־
$\alpha$
ושאר השורות לא השתנו. קיבלנו את 
$S(A)$.


\item
$S=(R_i+\alpha R_j\to R_i)$
עבור
$i\neq j$.
כאן
$$
,E = I_m + \alpha P(i,j)
$$
כאשר
$P(i,j)$ 
היא המטריצה שכל ערכיה הם
$0$
פרט ל-
$(P(i,j))_{ij}=1$.
נחשב:
$$
.EA = (I_m + \alpha P(i,j))A = A + \alpha (P(i,j)A)
$$
אבל
$$
(P(i,j)A)_{rk} = \sum_{l=1}^m (P(i,j))_{rl} a_{lk} = \begin{cases}
0, & r \neq i\\
a_{jk}, & r = i
\end{cases}
$$
ולכן
$$
(EA)_{rk} =
\begin{cases}
a_{rk}, & r \neq i\\
a_{ik} + \alpha a_{jk}, & r = i
\end{cases}
$$
כלומר השורה
$R_i$
הוחלפה ב־
$R_i + \alpha R_j$.
שוב קיבלנו את
$S(A)$.


\end{enumerate}
\end{proof}

\begin{corollary}\label{cor:left}
\addcorollarytolist{הקשר בין מטריצות שקולות שורה במונחים של כפל}
יהיו
$A,A'\in\M_{m\times n}(\F)$
שקולות שורה. אז קיימת 
$B\in\M_{m\times m}(\F)$
כך ש-
$A'=BA$.
\end{corollary}

\begin{proof}
לפי הנתון קיימות פש"אות 
$S_1,S_2,...,S_k$
כך שמתקיים
$
,A'=S_k(...(S_1(A))...)
$
כלומר 
$A'$
מתקבלת
מ-
$A$
בעקבות הפעלת
$S_1$,
אחריה 
$S_2$
וכן הלאה עד
$S_k$.

לפי טענה
\ref{prop:elementary}, 
ניתן לפתוח את הסוגריים בהדרגה כאשר בכל שלב הפעולה הרלוונטית הופכת לכפל במטריצה האלמנטרית המתאימה לה:

$$
\begin{aligned}
A'&=S_k(I_m)S_{k-1}(...(S_1(A))...)=S_k(I_m)S_{k-1}(I_m)S_{k-2}(...(S_1(A))...)=...\\
&=S_k(I_m)S_{k-1}(I_m)S_{k-2}(I_m)\cdots S_1(I_m)A
\end{aligned}
$$
נסמן
$B=S_k(I_m)S_{k-1}(I_m)S_{k-2}(I_m)\cdots S_1(I_m)$.
זו אכן מטריצה מסדר 
$m\times m$
וקיבלנו
$A'=BA$.
\end{proof}

\begin{example}
נסתכל על המטריצה

$$
.A =
\begin{pmatrix}
1 & 2 & 3 & 1 \\
2 & 4 & 5 & -2 \\
1 & 3 & 4 & 0
\end{pmatrix}
$$

נבצע דירוג קנוני, אבל הפעם נציין בכל שלב את המטריצה האלמנטרית המתאימה
$E_i$
ונעקוב אחר המכפלה המצטברת.

% ---------- שלב 1 ----------
שלב 1: $R_2 - 2R_1 \to R_2$
\begin{alignat*}{2}
E_1 &=
\begin{pmatrix}
1 & 0 & 0 \\
-2 & 1 & 0 \\
0 & 0 & 1
\end{pmatrix},
\qquad
E_1 A &=
\begin{pmatrix}
1 & 2 & 3 & 1 \\
0 & 0 & -1 & -4 \\
1 & 3 & 4 & 0
\end{pmatrix}
\end{alignat*}

% ---------- שלב 2 ----------
שלב 2: $R_3 - R_1 \to R_3$
\begin{alignat*}{2}
E_2 &=
\begin{pmatrix}
1 & 0 & 0 \\
0 & 1 & 0 \\
-1 & 0 & 1
\end{pmatrix},
\qquad
E_2 E_1 A &=
\begin{pmatrix}
1 & 2 & 3 & 1 \\
0 & 0 & -1 & -4 \\
0 & 1 & 1 & -1
\end{pmatrix}
\end{alignat*}

% ---------- שלב 3 ----------
שלב 3: $R_3 \leftrightarrow R_2$
\begin{alignat*}{2}
E_3 &=
\begin{pmatrix}
1 & 0 & 0 \\
0 & 0 & 1 \\
0 & 1 & 0
\end{pmatrix},
\qquad
E_3 E_2 E_1 A &=
\begin{pmatrix}
1 & 2 & 3 & 1 \\
0 & 1 & 1 & -1 \\
0 & 0 & -1 & -4
\end{pmatrix}
\end{alignat*}

% ---------- שלב 4 ----------
שלב 4: $-R_3 \to R_3$ 
\begin{alignat*}{2}
E_4 &=
\begin{pmatrix}
1 & 0 & 0 \\
0 & 1 & 0 \\
0 & 0 & -1
\end{pmatrix},
\qquad
E_4 E_3 E_2 E_1 A &=
\begin{pmatrix}
1 & 2 & 3 & 1 \\
0 & 1 & 1 & -1 \\
0 & 0 & 1 & 4
\end{pmatrix}
\end{alignat*}

% ---------- שלב 5 ----------
שלב 5: $R_2 - R_3 \to R_2$
\begin{alignat*}{2}
E_5 &=
\begin{pmatrix}
1 & 0 & 0 \\
0 & 1 & -1 \\
0 & 0 & 1
\end{pmatrix},
\qquad
E_5 E_4 E_3 E_2 E_1 A &=
\begin{pmatrix}
1 & 2 & 3 & 1 \\
0 & 1 & 0 & -5 \\
0 & 0 & 1 & 4
\end{pmatrix}
\end{alignat*}

% ---------- שלב 6 ----------
שלב 6: $R_1 - 3R_3 \to R_1$
\begin{alignat*}{2}
E_6 &=
\begin{pmatrix}
1 & 0 & -3 \\
0 & 1 & 0 \\
0 & 0 & 1
\end{pmatrix},
\qquad
E_6 E_5 E_4 E_3 E_2 E_1 A &=
\begin{pmatrix}
1 & 2 & 0 & -11 \\
0 & 1 & 0 & -5 \\
0 & 0 & 1 & 4
\end{pmatrix}
\end{alignat*}

% ---------- שלב 7 ----------
שלב 7: $R_1 - 2R_2 \to R_1$
\begin{alignat*}{2}
E_7 &=
\begin{pmatrix}
1 & -2 & 0 \\
0 &  1 & 0 \\
0 &  0 & 1
\end{pmatrix},
\qquad
E_7 E_6 E_5 E_4 E_3 E_2 E_1 A &=
\begin{pmatrix}
1 & 0 & 0 & -1 \\
0 & 1 & 0 & -5 \\
0 & 0 & 1 & 4
\end{pmatrix}
\end{alignat*}

המטריצה האחרונה שקיבלנו בצד ימין היא מדורגת קנונית:

$$
A'=
\begin{pmatrix}
1 & 0 & 0 & -1 \\
0 & 1 & 0 & -5 \\
0 & 0 & 1 & 4
\end{pmatrix}
$$

המכפלה הכוללת של המטריצות האלמנטריות היא
$E = E_7 E_6 E_5 E_4 E_3 E_2 E_1$,
ובמפורש:

$$
E =
\begin{pmatrix}
1 & 1 & -2 \\
-3 & 1 & 1 \\
2 & -1 & 0
\end{pmatrix}
$$

בדיקה:

$$
.EA =
\begin{pmatrix}
1 & 1 & -2 \\
-3 & 1 & 1 \\
2 & -1 & 0
\end{pmatrix}\begin{pmatrix}
1 & 2 & 3 & 1 \\
2 & 4 & 5 & -2 \\
1 & 3 & 4 & 0
\end{pmatrix}=\begin{pmatrix}
1 & 0 & 0 & -1 \\
0 & 1 & 0 & -5 \\
0 & 0 & 1 & 4
\end{pmatrix}
=A'
$$
\end{example}

\section{מטריצות הפיכות והקשר לממ"ל}


לפני שנמשיך, יש סימון נוח לוקטורים שיהיו שימושיים בקורס:

\begin{definition}
\adddeftolist{וקטורי היחידה}
נגדיר את וקטורי היחידה
$e_1,e_2,\cdots,e_n\in\F^n$
ע"י
$$
e_i=\begin{pmatrix}
0 \\
0 \\
\vdots \\
0 \\
1 \\
0\\
\vdots \\
0
\end{pmatrix}
$$
כאשר כל הקוארדינטות הן
$0$
חוץ מהקוארדינטה ה-
$i$
השווה ל-
$1$.
נציין שאלה וקטורי העמודה של 
$I_n$,
וגם וקטורי השורה עד כדי שחלוף (ההבדל בין שורה לעמודה).
\end{definition}

\begin{exercise}
לכל
$$
x=\begin{pmatrix}
    a_1 \\
    a_2 \\
    \vdots \\
    a_n
    \end{pmatrix}
$$
מתקיים
$$
.x=\sum_{j=1}^n a_je_j
$$
\end{exercise}

\begin{solution}
בסכום 
$
\sum_{j=1}^n a_je_j
$
רק המחובר ה-
$i$
תורם לקוארדינטה ה-
$i$,
כי כל השאר תורמים
$0$.
התרומה עבור
$j=i$
היא
$a_i\cdot 1=a_i$
כנדרש.
\end{solution}


לפני שנגדיר מטריצות הפיכות, נוכיח טענה שמראה את חשיבותה של מטריצת היחידה בהקשר של קיום ויחידות הפתרון לממ"ל כאשר מספר המשתנים שווה למספר המשוואות.

\begin{proposition}\label{prop:RREF}
\addpropositiontolist{הקשר בין דירוג קנוני, דרגה ויחידות הפתרון לממ"ל}
אם
$A'\in\M_{n\times n}(\F)$
מדורגת קנונית, אז התנאים הבאים שקולים:

\begin{enumerate}[label=(\alph*)]
\item 
$$A'=I_n$$


\item 
$$\rank(A')=n$$


\item 
לממ"ל
$A'x=\underline{0}$
יש פתרון יחיד.


\item 
לכל
$b'\in\F^n$
יש פתרון יחיד לממ"ל
$A'x=b'$.


\item 
לכל
$b'\in\F^n$
יש פתרון לממ"ל
$A'x=b'$.


\end{enumerate}
\end{proposition}

\begin{proof}
נוכיח כי
$(\text{א})$
שקול לכל אחד מהתנאים האחרים. תחילה נראה שהוא גורר את האחרים:

$(\text{ב}) \impliedby (\text{א})$:
למטריצה
$I_n$
יש דרגה 
$n$
כמספר האיברים המובילים שלה.

\noindent $(\text{ג}) \impliedby (\text{א})$:
ברור שיש פתרון יחיד, שהוא הפתרון הטריוויאלי.

\noindent $(\text{ד}) \impliedby (\text{א})$:
ברור שיש פתרון יחיד, שהוא
$x=b'$.

\noindent $(\text{ה}) \impliedby (\text{א})$:
כנ"ל.

כעת נוכיח את הכיוון ההפוך של כל גרירה.

\noindent $(\text{ב}) \implies (\text{א})$:
אם הדרגה היא 
$n$,
אז יש
$1$
מוביל בכל עמודה. כל שאר האיברים הם
$0$
,(כי הדירוג קנוני)
ולכן
$A'=I_n$.

\noindent $(\text{ג}) \implies (\text{א})$:
מיחידות הפתרון נובע 
$(\text{ב})$
ומכאן גם
$(\text{א})$.


\noindent $(\text{ד}) \implies (\text{א})$:
כנ"ל.


\noindent $(\text{ה}) \implies (\text{א})$:
אם תמיד קיים פתרון לממ"ל
$A'x=b$,
אז אין שורת אפסים ב-
$A'$
כי אחרת מתקבלת שורה סתירה ב-
$(A'|b')$
עבור
$$
.b'=\begin{pmatrix}
0\\
0\\
\vdots \\
0 \\
1
\end{pmatrix}
$$

לכן נובע
$(\text{ב})$
ומכאן
$(\text{א})$.

\end{proof}

כמסקנה מהטענה הקודמת למטריצות ריבועיות שהן מדורגות קנונית, יש לנו משפט למטריצות ריבועיות כלליות:

\begin{proposition}\label{prop:invertible}
לכל
$A\in\M_{n\times n}(\F)$
התנאים הבאים שקולים:
\begin{enumerate}[label=(\alph*)]
\item 
$A$
שקולת שורה ל-
$I_n$.


\item 
$\rank(A)=n$


\item 
לממ"ל
$Ax=\underline{0}$
יש פתרון יחיד.


\item 
לכל
$b\in\F^n$
יש פתרון יחיד לממ"ל
$Ax=b$.


\item 
לכל
$b\in\F^n$
יש פתרון לממ"ל
$Ax=b$.


\end{enumerate}
\end{proposition}

\begin{proof}
תהי 
$A\in\M_{n\times n}(\F)$.
היא שקולת שורה למטריצה מתאימה 
$A'$
שהיא מדורגת קנונית, ונשים לב כי כל אחד מהתנאים
$\text{(ה)},\text{(ד)},\text{(ג)},\text{(ב)}$
שקול לתנאי המקביל לו בטענה
\ref{prop:RREF}
עם
$A'$
במקום
$A$.
הסיבה לכך היא שמתקיים
$\rank(A)=\rank(A')$
לפי הגדרת דרגה, וגם קבוצת הפתרונות של הממ"ל
$Ax=b$
שווה לקבוצת הפתרונות של הממ"ל
$A'x=b'$.
לכן, כל אחד מהתנאים
$\text{(ה)},\text{(ד)},\text{(ג)},\text{(ב)}$
שקול לתנאי
$A'=I_n$,
שזה בדיוק
$\text{(א)}$.

\end{proof}

משפט זה מראה שיש משהו מיוחד לגבי מטריצות ריבועיות שהן שקולות שורה למטריצת יחידה. נראה שיש דרך חלופית להגדיר אותן, שקשורה לכפל מטריצות. בהתחלה ההגדרה תיראה שונה, ובסוף נוכיח שזו הגדרה שקולה.

\begin{definition}
\adddeftolist{מטריצה הפיכה}
מטריצה 
$A\in\M_{n\times n}(\F)$
נקראת הפיכה משמאל אם קיימת
$B\in\M_{n\times n}(\F)$
כך ש-
$BA=I_n$.
$A$
נקראת הפיכה מימין אם קיימת
$C\in\M_{n\times n}(\F)$
כך ש-
$AC=I_n$.
אם שני התנאים מתקיימים, 
$A$
נקראת הפיכה.
\end{definition}

נראה שהפיכות מימין והפיכות משמאל, הן למעשה שני תנאים שקולים (עבור מטריצה ריבועית). יותר מכך:

\begin{theorem}\label{thm:invertible}
לכל
$A\in\M_{n\times n}(\F)$
התנאים הבאים שקולים:
\begin{enumerate}[label=(\alph*)]
\item 
$A$
שקולת שורה ל-
$I_n$.


\item 
$\rank(A)=n$


\item 
לממ"ל
$Ax=\underline{0}$
יש פתרון יחיד.


\item 
לכל
$b\in\F^n$
יש פתרון יחיד לממ"ל
$Ax=b$.


\item 
לכל
$b\in\F^n$
יש פתרון לממ"ל
$Ax=b$.


\item 
$A$
הפיכה משמאל.

\item 
$A$
הפיכה מימין.

\item 
$A$
הפיכה.


\end{enumerate}
\end{theorem}

\begin{proof}
כבר הוכחנו שחמשת התנאים הראשונים שקולים. נראה כי
$$
(\text{ג}) \impliedby (\text{ו}) \impliedby (\text{א})
$$
וגם
$$
(\text{ז}) \iff (\text{ה})
$$
ומכאן נסיק שכל אחד מהתנאים
$
(\text{ח}) ,(\text{ז}) ,(\text{ו})
$
שקול לכל שאר התנאים.

\noindent $
:(\text{ו}) \impliedby (\text{א})
$
זה נובע ישירות ממסקנה
\ref{cor:left}
עם
$A'=I_n$.

\noindent $
:(\text{ג}) \impliedby (\text{ו})
$
נניח כי קיימת
$B\in\M_{n\times n}(\F)$
כך ש-
$BA=I_n$.
אז כדי לפתור את הממ"ל
$Ax=\underline{0}$
נכפיל את שני האגפים ב-
$B$
משמאל:
$$
B(Ax)=B\cdot\underline{0} \implies I_n x=\underline{0} \implies x=\underline{0}
$$
לכן יש פתרון יחיד.

\noindent $
:(\text{ז}) \impliedby (\text{ה})
$
המטרה היא למצוא
$C\in\M_{n\times n}(\F)$
כך ש-
$AC=I_n$.
נסמן את העמודות של 
$C$
ע"י
$v_1,v_2,...,v_n\in\F^n$,
ואז לפי טענה
\ref{prop:product-column}
מתקיים
$$
.AC=\begin{pmatrix}
| & | & & | \\
Av_1 & Av_2 & \cdots & Av_n \\
| & | & & |
\end{pmatrix}
$$

לכן
$$
.AC=I_n \iff \forall 1\leq j\leq n \quad Av_j=e_j
$$

אז קיבלנו
$n$
ממ"ליות,
ולכל אחת מהן יש פתרון
$v_j$
לפי 
$(\text{ה})$.
אז המטריצה המתאימה 
$$
C=\begin{pmatrix}
| & | & & | \\
v_1 & v_2 & \cdots & v_n \\
| & | & & |
\end{pmatrix}
$$
מקיימת
$AC=I_n$
כנדרש.

\noindent $
:(\text{ה}) \impliedby (\text{ז})
$
נניח כי קיימת
$C\in\M_{n\times n}(\F)$
כך ש-
$AC=I_n$.
יהי 
$$
b=\begin{pmatrix}
\beta_1 \\
\beta_2 \\
\vdots \\
\beta_n 
\end{pmatrix}
$$

ונראה שקיים פתרון לממ"ל
$Ax=b$.
בסימונים הקודמים לעמודות
$C$,
שוב מתקיים לכל 
$1\leq j\leq n$
$$
.Av_j=e_j
$$

אז ניקח
$$
x=\begin{pmatrix}
x_1 \\
x_2 \\
\vdots \\
x_n
\end{pmatrix}
$$
עם
$
.x=\sum_{j=1}^n \beta_j v_j
$
לכן נקבל
$$
.Ax=A(\sum_{j=1}^n \beta_j v_j)=\sum_{j=1}^n \beta_j Av_j=\sum_{j=1}^n \beta_j e_j=b
$$

ניתן להכניס את
$A$
לתוך הסכימה לפי טענה 
\ref{prop:matrix-vector}.
\end{proof}

לכאורה, הפיכות מימין והפיכות משמאל מתייחסות לשתי מטריצות
$B,C$
שיכולות להיות שונות, אבל ראינו כי הפיכות משמאל שקולה להפיכות מימין. אז אמור להיות קשר בין המטריצות, ואכן:

\begin{proposition}\label{prop:inverses}
\addpropositiontolist{הפיכות מימין והפיכות משמאל}
יהיו 
$A,B,C\in\M_{n\times n}(\F)$ 
כך ש-
$BA=I_n=AC$.
אז
$B=C$.
\end{proposition}

\begin{proof}
לפי התכונות של 
$I_n$
וחוק הקיבוץ לכפל מטריצות,
נובע כי
$$
.B=BI_n=B(AC)=(BA)C=I_nC=C
$$
\end{proof}

\begin{corollary}
יהיו 
$A,B\in\M_{n\times n}(\F)$.
אז מתקיים
$$
.AB=I_n \iff BA=I_n
$$
\end{corollary}

\begin{proof}
אם
$BA=I_n$,
אז 
$A$
הפיכה משמאל. לפי משפט
\ref{thm:invertible}
היא גם הפיכה מימין, ולכן קיימת
$C\in\M_{n\times n}(\F)$
כך ש-
$AC=I_n$.
אבל אז נובע מטענה
\ref{prop:inverses}
כי
$B=C$,
ולכן
$AB=I_n$.

הכיוון השני דומה, למשל משיקולי סימטריה (משחק תפקידים בין
$A$
ל-
$B$).
\end{proof}

האם בהינתן מטריצה הפיכה
$A$,
המטריצה
$B$
שמקיימת
$AB=I_n=BA$
נקבעת ביחידות ע"י
$A$? 
התשובה חיובית.
לפני שנתאר את החישוב, נוכיח את היחידות:

\begin{corollary}
\addpropositiontolist{יחידות של מטריצה הופכית}
יהיו
$A,B_1,B_2\in\M_{n\times n}(\F)$
כך ש-
$B_1A=B_2A=I_n$.
אז 
$B_1=B_2$.
\end{corollary}

\begin{proof}
$A$
הפיכה משמאל ולכן גם מימין לפי משפט
\ref{thm:invertible}.
כלומר קיימת
$C\in\M_{n\times n}(\F)$
כך ש-
$AC=I_n$.
לכן, נובע
$B_1=C=B_2$
מטענה
\ref{prop:inverses}.
\end{proof}

\begin{definition}
\adddeftolist{מטריצה הופכית}
תהי
$A\in\M_{n\times n}(\F)$
הפיכה. נסמן ב-
$A^{-1}$
את המטריצה שמקיימת
$$
.AA^{-1}=A^{-1}A=I_n
$$
נקרא לה המטריצה ההופכית ל-
$A$.
\end{definition}

\begin{remark}
נשים לב כי 
$A,A^{-1}$
מתחלפות. בהינתן 
$B$,
כדי לבדוק שמתקיים
$B=A^{-1}$
מספיק להראות
$AB=I_n$
או לחילופין
$BA=I_n$,
כי התנאים שקולים.
\end{remark}

\begin{exercise}\label{exer:inverse-product}
יהיו
$A,B\in\M_{n\times n}(\F)$
הפיכות ו-
$0\neq \alpha\in\F$.
אז מתקיים:
\begin{enumerate}[label=\alph*.]
\item 
$A^{-1}$
הפיכה ומתקיים
$$
.(A^{-1})^{-1}=A
$$

\item
$\alpha A$
הפיכה ומתקיים
$$
.(\alpha A)^{-1}=\frac{1}{\alpha}A^{-1}
$$


\item
$AB$
הפיכה ומתקיים
$$
.(AB)^{-1}=B^{-1}A^{-1}
$$


\item 
$A^t$
הפיכה ומתקיים
$$
.(A^t)^{-1}=(A^{-1})^t
$$


\end{enumerate}
\end{exercise}

\begin{solution}\
בכל סעיף נחשב את אחת משתי המכפלות ונראה שהיא שווה ל-
$I_n$.
\begin{enumerate}[label=\alph*.]
\item 
בהינתן
$AA^{-1}=A^{-1}A=I_n$
נובע כי גם
$A^{-1}$
הפיכה וההופכית שלה היא
$A$.

\item
מתקיים
$$
.(\frac{1}{\alpha}A^{-1})(\alpha A)=(\frac{1}{\alpha}\alpha)A^{-1}A=1\cdot I_n=I_n
$$
לכן
$\alpha A$
הפיכה ומתקיים
$
.(\alpha A)^{-1}=\frac{1}{\alpha}A^{-1}
$


\item
מתקיים
$$
.(AB)(B^{-1}A^{-1})=A(BB^{-1})A^{-1}=A(I_nA^{-1})=AA^{-1}=I_n
$$
לכן
$AB$
הפיכה ומתקיים
$
(AB)^{-1}=B^{-1}A^{-1}
$.


\item
מתקיים
$$
.A^t(A^{-1})^t=(A^{-1}A)^t=I_n^t=I_n
$$

לכן 
$A^t$
הפיכה ומתקיים
$
.(A^t)^{-1}=(A^{-1})^t
$


\end{enumerate}
\end{solution}

\begin{remark}
לא בהכרח מתקיים
$
(AB)^{-1}=A^{-1}B^{-1}
$
כי אין הכרח ש-
$A^{-1},B^{-1}$
יתחלפו.
\end{remark}

\begin{remark}
אם
$A$
הפיכה, אז הפתרון היחיד לממ"ל
$Ax=b$
נתון ע"י
$$x=A^{-1}b$$
לכל
$b\in\F^n$.
אכן, ניתן להכפיל את שני אגפי המשוואה ב-
$A^{-1}$
בהינתן שהיא קיימת. לעומת זאת, אי אפשר "לחלק" במטריצה שאינה הפיכה.
\end{remark}

\subsection{חישוב מטריצה הופכית}

הבנו שכל מטריצה
$A\in\M_{n\times n}(\F)$
הפיכה שקולת שורה ל-
$I_n$.
יותר מכך, לפי ההוכחה של מסקנה
\ref{cor:left}
יש התאמה בין 
$A^{-1}$
לבין הפש"אות שמובילות ל-
$I_n$
בדירוג של 
$A$.
נחדד את נקודה זו כמסקנה נוספת:

\begin{corollary}\label{cor:elementary-product}\
\addcorollarytolist{הקשר בין מטריצות הפיכות למטריצות אלמנטריות}
\begin{enumerate}[label=\alph*.]
\item
תהי 
$A\in\M_{n\times n}(\F)$
הפיכה. אז מתקיים
$$
.A^{-1}=E_kE_{k-1}\cdots E_1
$$
כאשר
$E_1,E_2,...,E_k$
הן המטריצות האלמנטריות המתאימות לפש"אות
$S_1,S_2,...,S_k$
בדירוג הקנוני של 
$A$.


\item 
$A$
היא מכפלה של מטריצות אלמנטריות:
$$
.A=E_1^{-1}E_2^{-1}\cdots E_k^{-1}
$$


\end{enumerate}
\end{corollary}

\begin{proof}\
\begin{enumerate}[label=\alph*.]
\item
זהו מקרה פרטי של מסקנה
\ref{cor:left}
עבור
$A'=I_n$.
כאן
$B=A^{-1}$.


\item 
לפי תרגיל
\ref{exer:inverse-product}
הפעולה של הפיכת מטריצה הופכת את סדר המטריצות במכפלה. אפשר להבין זאת מחישוב הדרגתי:

$$
\begin{aligned}
A^{-1}&=E_kE_{k-1}\cdots E_1 \\
\implies A&=((E_kE_{k-1}\cdots E_2) E_1)^{-1}=E_1^{-1}((E_kE_{k-1}\cdots E_3)E_2)^{-1}\\
&=E_1^{-1}E_2^{-1}((E_kE_{k-1}\cdots E_4)E_3)^{-1}=...=E_1^{-1}E_2^{-1}\cdots E_k^{-1}
\end{aligned}
$$

לחילופין, אפשר לוודא שמתקיים
$$
.(E_1^{-1}E_2^{-1}\cdots E_k^{-1})(E_kE_{k-1}\cdots E_1)=I_n
$$

נדגיש שלכל מטריצה אלמנטרית, גם ההופכית שלה היא מטריצה אלמנטרית המתאימה לפש"א ההפוכה.

\end{enumerate}
\end{proof}

יש דרך נוחה ומסודרת לחשב את המטריצות האלמנטריות ואת מכפלתן שמובילה למטריצה ההופכית. הרעיון הוא לדרג את 
$I_n$
לצד
$A$,
כמו במטריצת מקדמים מורחבת - אבל עם מספר עמודות מימין לקו המפריד. קובעים את הפש"אות לפי 
$A$,
אבל מבצעים אותן גם על
$I_n$ 
(כלומר מכפילים גם אותה במטריצות האלמנטריות המתאימות):

$$
(A\mid I_n)\rightarrow (E_1A\mid E_1)\rightarrow (E_2E_1A\mid E_2E_1)\rightarrow \cdots \rightarrow(E_k\cdots E_1A \mid E_k\cdots E_1)
$$

המטריצה האחרונה בצד ימין היא למעשה
$(I_n \mid A^{-1})$,
כלומר המטריצה ההופכית מופיעה מימין ל-
$I_n$
בסוף הדירוג הקנוני של מטריצה הפיכה.

% Insert link to Inverse.html here.

\begin{example}
\begin{enumerate}[label=\alph*.]
\item
ניקח את
$$
A=\begin{pmatrix}
1 & 2 & 1 \\
0 & 1 & 1 \\
2 & 3 & 4 
\end{pmatrix}
$$
ונחשב את ההופכית ע"י דירוג קנוני של 
$(A \mid I_3)$:

\begin{align*}
(A\mid I_3)
&=\left(
\begin{array}{ccc|ccc}
1 & 2 & 1 & 1 & 0 & 0\\
0 & 1 & 1 & 0 & 1 & 0\\
2 & 3 & 4 & 0 & 0 & 1
\end{array}
\right)
\\[4pt]
&\xrightarrow{R_3-2R_1\to R_3}
\left(
\begin{array}{ccc|ccc}
1 & 2 & 1 & 1 & 0 & 0\\
0 & 1 & 1 & 0 & 1 & 0\\
0 & -1 & 2 & -2 & 0 & 1
\end{array}
\right)
\\[4pt]\displaybreak[2]
&\xrightarrow{R_2\leftrightarrow R_3}
\left(
\begin{array}{ccc|ccc}
1 & 2 & 1 & 1 & 0 & 0\\
0 & -1 & 2 & -2 & 0 & 1\\
0 & 1 & 1 & 0 & 1 & 0
\end{array}
\right)
\\[4pt]
&\xrightarrow{-R_2\to R_2}
\left(
\begin{array}{ccc|ccc}
1 & 2 & 1 & 1 & 0 & 0\\
0 & 1 & -2 & 2 & 0 & -1\\
0 & 1 & 1 & 0 & 1 & 0
\end{array}
\right)
\\[4pt]
&\xrightarrow{R_1-2R_2\to R_1}
\left(
\begin{array}{ccc|ccc}
1 & 0 & 5 & -3 & 0 & 2\\
0 & 1 & -2 & 2 & 0 & -1\\
0 & 1 & 1 & 0 & 1 & 0
\end{array}
\right)
\\[4pt]
&\xrightarrow{R_3-R_2\to R_3}
\left(
\begin{array}{ccc|ccc}
1 & 0 & 5 & -3 & 0 & 2\\
0 & 1 & -2 & 2 & 0 & -1\\
0 & 0 & 3 & -2 & 1 & 1
\end{array}
\right)
\\[4pt]
&\xrightarrow{\tfrac13 R_3\to R_3}
\left(
\begin{array}{ccc|ccc}
1 & 0 & 5 & -3 & 0 & 2\\
0 & 1 & -2 & 2 & 0 & -1\\
0 & 0 & 1 & -\tfrac23 & \tfrac13 & \tfrac13
\end{array}
\right)
\\[4pt]
&\xrightarrow{R_1-5R_3\to R_1}
\left(
\begin{array}{ccc|ccc}
1 & 0 & 0 & \tfrac13 & -\tfrac53 & \tfrac13\\
0 & 1 & -2 & 2 & 0 & -1\\
0 & 0 & 1 & -\tfrac23 & \tfrac13 & \tfrac13
\end{array}
\right)
\\[4pt]
&\xrightarrow{R_2+2R_3\to R_2}
\left(
\begin{array}{ccc|ccc}
1 & 0 & 0 & \tfrac13 & -\tfrac53 & \tfrac13\\
0 & 1 & 0 & \tfrac23 & \tfrac23 & -\tfrac13\\
0 & 0 & 1 & -\tfrac23 & \tfrac13 & \tfrac13
\end{array}
\right)
\end{align*}

אז קיבלנו
$$
.A^{-1}=
\begin{pmatrix}
\frac{1}{3} & -\frac{5}{3} & \frac{1}{3}\\
\frac{2}{3} & \frac{2}{3} & -\frac{1}{3}\\
-\frac{2}{3} & \frac{1}{3} & \frac{1}{3}
\end{pmatrix}
$$


\item 
ניקח את
$$
B=\begin{pmatrix}
1+i & 2 \\
3 & 4-i
\end{pmatrix}
$$
ונחשב את ההופכית ע"י דירוג קנוני של 
$(B \mid I_2)$:

$$
\begin{aligned}
(B\mid I_2)
&=\left(
\begin{array}{cc|cc}
1+i & 2 & 1 & 0 \\
3 & 4-i & 0 & 1
\end{array}
\right)
\\[4pt]
&\xrightarrow{R_2-\tfrac{3}{1+i}R_1\to R_2}
\left(
\begin{array}{cc|cc}
1+i & 2 & 1 & 0\\
0 & 4-i - \tfrac{6}{1+i} & -\tfrac{3}{1+i} & 1
\end{array}
\right)
\\[4pt]
&\xrightarrow{\tfrac{1}{1+i}R_1\to R_1}
\left(
\begin{array}{cc|cc}
1 & \tfrac{2}{1+i} & \tfrac{1}{1+i} & 0\\
0 & \tfrac{-1+3i}{1+i} & -\tfrac{3}{1+i} & 1
\end{array}
\right)
\\[4pt]
&\xrightarrow{\tfrac{1+i}{-1+3i}R_2\to R_2}
\left(
\begin{array}{cc|cc}
1 & 1-i & \tfrac{1-i}{2} & 0\\
0 & 1 & \tfrac{3}{1-3i} & \tfrac{1+i}{-1+3i}
\end{array}
\right)
\\[4pt]
&\xrightarrow{R_1-(1-i)R_2\to R_1}
\left(
\begin{array}{cc|cc}
1 & 0 & \tfrac{1-i}{2}-\tfrac{(3+9i)(1-i)}{10} & -\tfrac{(1-i)(1-2i)}{5} \\
0 & 1 & \tfrac{3+9i}{10} & \tfrac{1-2i}{5}
\end{array}
\right)
\end{aligned}
$$

נפשט את המטריצה האחרונה שקיבלנו בצד ימין:

$$
B^{-1}=\begin{pmatrix}
-\tfrac{7+11i}{10} & \tfrac{1+3i}{5} \\
\tfrac{3+9i}{10} & \tfrac{1-2i}{5}
\end{pmatrix}
$$


\end{enumerate}
\end{example}

\begin{exercise}\
\begin{enumerate}[label=\alph*.]
\item
חשבו את
$A^{-1}$
עבור
$$
.A=\begin{pmatrix}
1 & 2 \\
3 & 4
\end{pmatrix}
$$

\item 
עבור
$A$
הקודמת, מצאו
$B\in\M_{2\times 2}(\R)$
עבורה מתקיים
$
AB=C
$
כאשר
$$
.C=\begin{pmatrix}
3 & -5 \\
-9 & 1
\end{pmatrix}
$$


\end{enumerate}
\end{exercise}

\begin{solution}\
\begin{enumerate}[label=\alph*.]
\item
נדרג את
$(A|I_2)$:

\begin{align*}
\left(A \mid I_2\right)
&=
\left(
\begin{array}{cc|cc}
1 & 2 & 1 & 0\\
3 & 4 & 0 & 1
\end{array}
\right)
\\[4pt]
&\xrightarrow{R_2-3R_1\to R_2}
\left(
\begin{array}{cc|cc}
1 & 2 & 1 & 0\\
0 & -2 & -3 & 1
\end{array}
\right)
\\[4pt]
&\xrightarrow{-\tfrac12 R_2\to R_2}
\left(
\begin{array}{cc|cc}
1 & 2 & 1 & 0\\
0 & 1 & \tfrac32 & -\tfrac12
\end{array}
\right)
\\[4pt]
&\xrightarrow{R_1-2R_2\to R_1}
\left(
\begin{array}{cc|cc}
1 & 0 & -2 & 1\\
0 & 1 & \tfrac32 & -\tfrac12
\end{array}
\right)
\end{align*}

אז קיבלנו

$$
.A^{-1}=\begin{pmatrix}
-2 & 1\\
\tfrac32 & -\tfrac12
\end{pmatrix}
$$

\item 
נכפיל משמאל ב-
$A^{-1}$
את שני אגפי המשוואה
$AB=C$
ונקבל:

$$
.B=A^{-1}C=\begin{pmatrix}
-2 & 1\\
\tfrac32 & -\tfrac12
\end{pmatrix}\begin{pmatrix}
3 & -5 \\
-9 & 1
\end{pmatrix}=\begin{pmatrix}
-15 & 11 \\
9 & -8 
\end{pmatrix}
$$


\end{enumerate}
\end{solution}

\clearpage
\section{סיכום}

\noindent\textbf{הגדרה 6.1.}

נגדיר את קבוצת כל המטריצות מסדר 
$m\times n$
מעל
$\F$
(כאשר
$\F=\R$
או
$\F=\C$
)
ע"י 

$$
.\M_{m\times n}(\F)=\Set{\begin{pmatrix}a_{11} & a_{12} & \cdots & a_{1n}\\
a_{21} & a_{22} & \cdots & a_{2n}\\
\vdots & \vdots & \ddots & \vdots\\
a_{m1} & a_{m2} & \cdots & a_{mn}
\end{pmatrix}|\forall 1\leq i\leq m \quad \forall 1\leq j\leq n \quad a_{ij}\in \F}
$$
\bigskip
\bigskip

\noindent\textbf{הגדרה 6.2.}

עבור וקטורים
$ה=(a_1,...,a_n)\in\F^n,w=(b_1,...,b_n)\in\F^n$
וסקלר
$c\in\F$
נגדיר את הפעולות הבאות:
\begin{enumerate}[label=\alph*.]
\item 
פעולת החיבור:
$$
v+w=(a_1+b_1,...,a_n+b_n)
$$


\item
פעולת כפל בסקלר:
$$
c \cdot v=(c\cdot a_1,\dots,c\cdot a_n)
$$


\end{enumerate}
\bigskip
\bigskip

\noindent\textbf{הגדרה 6.3.}

עבור
$A,B\in\M_{m\times n}(\F)$
וסקלר
$c\in\F$
נגדיר את הפעולות הבאות:

\begin{enumerate}[label=\alph*.]
\item 
פעולת החיבור:
$$
\begin{aligned}
A+B&=\begin{pmatrix}a_{11} & \cdots & a_{1n}\\
a_{21} & \cdots & a_{2n}\\
\vdots  & \ddots & \vdots\\
a_{m1}  & \cdots & a_{mn}
\end{pmatrix}+\begin{pmatrix}b_{11} & \cdots & b_{1n}\\
b_{21} & \cdots & b_{2n}\\
\vdots & \ddots & \vdots\\
b_{m1} & \cdots & b_{mn}
\end{pmatrix}\\
&=\begin{pmatrix}a_{11}+b_{11} & \cdots & a_{1n}+b_{1n} \\ a_{21}+b_{21} & \cdots & a_{2n}+b_{2n}\\
\vdots & \ddots & \vdots\\
a_{m1}+b_{m1} & \cdots & a_{mn}+b_{mn}
\end{pmatrix}
\end{aligned}
$$


\item
פעולת כפל בסקלר:
$$
c A=\begin{pmatrix}ca_{11} & ca_{12} & \cdots & ca_{1n}\\
ca_{21} & ca_{22} & \cdots & ca_{2n}\\
\vdots & \vdots & \ddots & \vdots\\
ca_{m1} & ca_{m2} & \cdots & ca_{mn}
\end{pmatrix}
$$

בפרט, נגדיר
$-A=(-1)\cdot A$ 
ובהתאמה 
$B-A=B+(-1)\cdot A$.


\end{enumerate}
\bigskip
\bigskip

\noindent\textbf{הגדרה 6.4.}

יש דרך מקוצרת לכתיבת חיבור. במקום לכתוב את החיבור בכל רכיב, נסתכל על האיבר הכללי של כל מטריצה. נסמן את האיבר בשורה ה-
$i$
ובעמודה ה-
$j$
ע"י
$(P)_{ij}$
לכל
$P\in\M_{m\times n}(\F)$. 
אז למעשה הגדרנו
לכל
$1\leq i \leq m$
ולכל
$1\leq j\leq n$:
$$(A+B)_{ij}=(A)_{ij}+(B)_{ij}$$

$$(c A)_{ij}=c(A)_{ij}$$
\bigskip
\bigskip

\noindent\textbf{הגדרה 6.5.}

לכל 
$m,n\in\N$
נגדיר את מטריצת האפס מסדר
$m\times n$
כמטריצה שכל איבריה הם
$0$:
$$
.\0_{m\times n}=\begin{pmatrix}0 &0 &\cdots &0 \\
0 &0 &\cdots &0 \\
\vdots & \vdots &\ddots &\vdots\\
0 &0 &\cdots &0
\end{pmatrix}
$$
\bigskip
\bigskip

\noindent\textbf{טענה 6.6.}

יהיו 
$A,B,C\in\M_{m\times n}(\F)$
ו-
$\alpha,\beta\in\F$.
אז מתקיימים הכללים הבאים:

\begin{enumerate}[label=\alph*.]
\item
חוק החילוף (קומוטטיביות): 
$$A+B=B+A$$


\item 
חוק הקיבוץ (אסוציאטיביות):
$$(A+B)+C=A+(B+C)$$


\item 
חוק הקיבוץ לכפל בסקלר:
$$(\alpha\beta)A=\alpha(\beta A)$$


\item 
חוק הפילוג (דיסטריביוטיביות):
$$\alpha(A+B)=\alpha A+\alpha B$$


\item
ניטרליות היחידה:
$$1\cdot A=A$$


\item
$$0\cdot A=\0_{m\times n}$$


\item 
ניטרליות מטריצת האפס:
$$A+\0_{m\times n}=A$$


\item 
$$A-A=0_{m\times n}$$


\end{enumerate}
\bigskip
\bigskip

\noindent\textbf{הגדרה 6.7.}

בהינתן
$a_1,a_2,...,a_n\in\F$
נסמן את סכומם בעזרת כתיבה עם
$\Sigma$
ואינדקס
$l$.

$$
.\sum_{l=1}^n a_l=a_1+a_2+a_3+\cdots+a_n
$$

האינדקס
$l$
רץ על פני הקבוצה
$\Set{1,2,3,\cdots,n}$
וכל הצבה תורמת איבר לסכום.
\bigskip
\bigskip

\noindent\textbf{הגדרה 6.8.}

בהינתן
$A\in\M_{m\times n}(\F)$
ו-
$$
,x=\begin{pmatrix}
x_1\\
x_2\\
\vdots \\
x_n
\end{pmatrix}\in\F^n
$$
נגדיר את המכפלה
$$
Ax=\begin{pmatrix}
y_1\\
y_2\\
\vdots \\
y_m
\end{pmatrix}\in\F^m
$$
עם קוארדינטות הנתונות ע"י
$$
y_i=\sum_{j=1}^n (A)_{ij}x_j
$$
לכל
$1\leq i\leq m$.
נסמן
$y_i=(Ax)_i$ 
עבור הקוארדינטות.
\bigskip
\bigskip

\noindent\textbf{טענה 6.9.}

יהיו
$A,B\in\M_{m\times n}(\F)$,
$v,w\in\F^n$
ו-
$\alpha\in\F$.
\begin{enumerate}[label=\alph*.]
\item
$$
A(v+w)=Av+Aw
$$


\item 
$$
(A+B)v=Av+Bv
$$


\item 
$$
A(\alpha v)=\alpha(Av)
$$


\end{enumerate}
\bigskip
\bigskip

\noindent\textbf{טענה 6.10.}

תהי
$Ax=\underline{0}$
ממ"ל הומוגנית.
\begin{enumerate}[label=\alph*.]
\item
קבוצת הפתרונות סגורה לחיבור, כלומר לכל שני פתרונות
$v,w\in\F^n$
גם
$v+w$
הוא פתרון.


\item
קבוצת הפתרונות סגורה לכפל בסקלר, כלומר לכל פתרון
$v\in\F^n$
ולכל
$\alpha\in\F$
גם
$\alpha v$
הוא פתרון.


\end{enumerate}
\bigskip
\bigskip

\noindent\textbf{טענה 6.11.}

תהי
$Ax=b$
ממ"ל לא הומוגנית. נניח שקיים לה פתרון (לא בהכרח יחיד)
$v_0\in\F^n$.
אז מתקיים הקשר הבא בין קבוצת הפתרונות של הממ"ל הלא הומוגנית לקבוצת הפתרונות של הממ"ל ההומוגנית
$Ax=\underline{0}$:

$$
.\Set{x\in\F^n|Ax=b}=\Set{v_0+y|y\in\F^n,Ay=\underline{0}}
$$
\bigskip
\bigskip

\noindent\textbf{הגדרה 6.12.}

בהינתן
$A\in\M_{m\times k}(\F),B\in\M_{k\times n}(\F)$
נגדיר את מטריצת הכפל
$AB$
ע"י
$$
(AB)_{ij}=\sum_{l=1}^k (A)_{il}(B)_{lj}=(A)_{i1}(B)_{1j}+(A)_{i2}(B)_{2j}+...+(A)_{ik}(B)_{kj}
$$
לכל
$1\leq i\leq m, 1\leq j\leq n$.
\bigskip
\bigskip

\noindent\textbf{טענה 6.13.}

יהיו
$A\in\M_{m\times k}(\F), B\in\M_{k\times n}(\F)$.
נניח שהעמודות של
$B$
הם הוקטורים
$v_1,v_2,...,v_n\in\F^k$.
אז מתקיים

$$
AB=A\begin{pmatrix}
| & | & & | \\
v_1 & v_2 & \cdots & v_n \\
| & | & & |
\end{pmatrix}=\begin{pmatrix}
| & | & & | \\
Av_1 & Av_2 & \cdots & Av_n \\
| & | & & |
\end{pmatrix}
$$

כלומר התוצאה של כפל ב-
$A$
משמאל הוא שכל וקטור עמודה 
מוכפל ב-
$A$
לחוד.
\bigskip
\bigskip

\noindent\textbf{טענה 6.14.}

יהיו
$A_1,A_2\in\M_{m\times n}(\F)$,
$B_1,B_2\in\M_{n\times p}(\F)$,
$C\in\M_{p\times q}(\F)$
ו-
$\alpha\in\F$.
אז מתקיים

\begin{enumerate}[label=\alph*.]
\item 
$$
(A_1+A_2)B_1=A_1B_1+A_2B_1
$$


\item 
$$
A_1(B_1+B_2)=A_1B_1+A_1B_2
$$


\item 
$$
A_1(B_1C)=(A_1B_1)C
$$


\item 
$$
(\alpha A_1)B_1=\alpha(A_1B_1)=A_1(\alpha B_1)
$$


\item 
$$
\0_{m\times n} B_1=\0_{m\times p}=A_1\0_{n\times p}
$$


\end{enumerate}
\bigskip
\bigskip

\noindent\textbf{הגדרה 6.15.}

מטריצת היחידה מסדר 
$n\times n$
(אפשר גם לומר מסדר $n$)
מוגדרת להיות
$$
.I_n=I_{n\times n}=
\begin{pmatrix}
1 &0 &\cdots &0\\
0 &1 &\cdots &0\\
\vdots &\vdots &\ddots &\vdots\\
0 &0 &\cdots &1
\end{pmatrix}
$$

זוהי המטריצה שבה כל איבר על האלכסון הראשי הוא
$1$,
וכל שאר האיברים הם
$0$. 
או באופן מתמטי:

$$
\begin{aligned}
\forall 1\leq i\leq n \quad (I_n)_{ii}=1 \\
\forall 1\leq i,j\leq n  \quad i\neq j \implies &(I_n)_{ij}=0
\end{aligned}
$$
\bigskip
\bigskip

\noindent\textbf{טענה 6.16.}

לכל 
$A\in\M_{m\times n}(\F)$
מתקיים
$AI_n=A=I_m A$.
\bigskip
\bigskip

\noindent\textbf{הגדרה 6.17.}

תהי
$A\in\M_{m\times n}(\F)$. 
נגדיר את המטריצה המשוחלפת
$A^t\in\M_{n\times m}(\F)$
ע"י
$$
(A^t)_{ij}=A_{ji}
$$
לכל
$1\leq i\leq n,1\leq j\leq m$.
\bigskip
\bigskip

\noindent\textbf{טענה 6.18.}

לכל 
$A,B\in\M_{m\times n}(\F), C\in\M_{n\times k}(\F)$
ולכל
$\alpha\in\F$
מתקיים:
\begin{enumerate}[label=\alph*.]
\item 
$$
(A^t)^t=A
$$


\item 
$$
(A+B)^t=A^t+B^t
$$

\item 
$$
(\alpha A)^t=\alpha A^t
$$


\item 
$$
(AC)^t=C^tA^t
$$


\end{enumerate}
\bigskip
\bigskip

\noindent\textbf{הגדרה 6.19.}

מטריצה 
$A$
נקראת ריבועית אם מספר השורות שלה שווה למספר העמודות, כלומר אם קיים 
$n\in\N$
עבורו
$A\in\M_{n\times n}(\F)$.
\bigskip
\bigskip

\noindent\textbf{הגדרה 6.20.}

תהי 
$A\in\M_{n\times n}(\F)$.
אז נגדיר 

$$
A^0=I_n, \quad A^1=A, \quad A^2=A\cdot A, \quad A^3=A\cdot A\cdot A
$$

ולכל
$k\in\N$

$$
.A^k = \underbrace{A \cdot A \cdot \ldots \cdot A}_{\text{פעמים} \ k}
$$

עבור פולינום
$$
p(x)=a_mx^m+a_{m-1}x^{m-1}+...+a_1x+a_0
$$
נגדיר את ההצבה של
$A$
בו ע"י
$$
.p(A)=a_mA^m+a_{m-1}A^{m-1}+...+a_1A+a_0I_n
$$
\bigskip
\bigskip

\noindent\textbf{הגדרה 6.21.}

תהי 
$A\in\M_{n\times n}(\F)$
מטריצה ריבועית.

\begin{enumerate}[label=\alph*.]
\item
$A$
נקראת סימטרית אם מתקיים
$A^t=A$,
או באופן שקול:
$$
\forall 1\leq i,j\leq n \quad (A)_{ji}=(A)_{ij}
$$


\item 
$A$
נקראת אנטי-סימטרית אם מתקיים
$A^t=-A$,
או באופן שקול:
$$
\forall 1\leq i,j\leq n \quad (A)_{ji}=-(A)_{ij}
$$

\end{enumerate}
\bigskip
\bigskip

\noindent\textbf{טענה 6.22.}

תהי
$P\in\M_{n\times n}(\F)$
מטריצה ריבועית.
\begin{enumerate}[label=\alph*.]
\item
אם
$P$
אנטי-סימטרית, אז כל איבריה על האלכסון הראשי שווים ל-
$0$.


\item 
אם
$P$
סימטרית וגם אנטי-סימטרית, אז
$P=\0_{n\times n}$.


\item 
לכל
$B\in\M_{n\times n}(\F)$
קיימות
$S\in\M_{n\times n}(\F)$
סימטרית ו-
$A\in\M_{n\times n}(\F)$
אנטי-סימטרית כך שמתקיים
$B=S+A$.


\end{enumerate}
\bigskip
\bigskip

\noindent\textbf{הגדרה 6.23.}

מטריצה 
$A\in\M_{n\times n}(\F)$
נקראת:

\begin{enumerate}[label=\alph*.]
\item 
אלכסונית אם לכל
$i\neq j$
מתקיים
$(A)_{ij}=0$,
כלומר כל האיברים מחוץ לאלכסון הראשי שווים ל-
$0$.


\item
משולשית עליונה אם לכל
$i>j$
מתקיים
$(A)_{ij}=0$,
כלומר כל האיברים מתחת לאלכסון הראשי שווים ל-
$0$.


\item
משולשית תחתונה אם לכל
$i<j$
מתקיים
$(A)_{ij}=0$,
כלומר כל האיברים מעל לאלכסון הראשי שווים ל-
$0$.


\end{enumerate}
\bigskip
\bigskip

\noindent\textbf{טענה 6.24.}

יהיו
$A,B\in\M_{n\times n}(\F)$.
\begin{enumerate}[label=\alph*.]
\item
אם
$A,B$
משולשיות עליונות, אז
$AB$
גם משולשית עליונה.


\item
אם
$A,B$
משולשיות תחתונות, אז
$AB$
גם משולשית תחתונה.


\item
אם
$A,B$
אלכסוניות, אז
$AB$
גם אלכסונית.


\end{enumerate}
\bigskip
\bigskip

\noindent\textbf{הגדרה 6.25.}

מטריצה 
$A\in\M_{n\times n}(\F)$
נקראת סקלרית אם קיים
$\alpha\in\F$
כך ש-
$A=\alpha I_n$.
\bigskip
\bigskip

\noindent\textbf{הגדרה 6.26.}

מטריצה
$A\in\M_{n\times n}(\F)$
נקראת אלמנטרית אם היא מתקבלת ממטריצת היחידה
$I_n$
ע"י פש"א אחת.

עבור פש"א 
$S$,
נסמן ב-
$S(I_n)$
את המטריצה האלמנטרית המתאימה שמתקבלת מ-
$I_n$
ע"י 
$S$.
באופן כללי, לכל מטריצה
$A$
נסמן ב-
$S(A)$
את המטריצה שמתקבלת מ-
$A$
ע"י
$S$.
\bigskip
\bigskip

\noindent\textbf{טענה 6.27.}

יהיו
$A\in\M_{m\times n}(\F)$
ו-
$S$
פש"א. אז מתקיים
$S(I_m)A=S(A)$.
\bigskip
\bigskip

\noindent\textbf{מסקנה 6.28.}

יהיו
$A,A'\in\M_{m\times n}(\F)$
שקולות שורה. אז קיימת 
$B\in\M_{m\times m}(\F)$
כך ש-
$A'=BA$.
\bigskip
\bigskip

\noindent\textbf{הגדרה 6.29.}

נגדיר את וקטורי היחידה
$e_1,e_2,\cdots,e_n\in\F^n$
ע"י
$$
e_i=\begin{pmatrix}
0 \\
0 \\
\vdots \\
0 \\
1 \\
0\\
\vdots \\
0
\end{pmatrix}
$$
כאשר כל הקוארדינטות הן
$0$
חוץ מהקוארדינטה ה-
$i$
השווה ל-
$1$.
נציין שאלה וקטורי העמודה של 
$I_n$,
וגם וקטורי השורה עד כדי שחלוף (ההבדל בין שורה לעמודה).
\bigskip
\bigskip

\noindent\textbf{טענה 6.30.}

אם
$A'\in\M_{n\times n}(\F)$
מדורגת קנונית, אז התנאים הבאים שקולים:

\begin{enumerate}[label=(\alph*)]
\item 
$$A'=I_n$$


\item 
$$\rank(A')=n$$


\item 
לממ"ל
$A'x=\underline{0}$
יש פתרון יחיד.


\item 
לכל
$b'\in\F^n$
יש פתרון יחיד לממ"ל
$A'x=b'$.


\item 
לכל
$b'\in\F^n$
יש פתרון לממ"ל
$A'x=b'$.


\end{enumerate}
\bigskip
\bigskip

\noindent\textbf{טענה 6.31.}

לכל
$A\in\M_{n\times n}(\F)$
התנאים הבאים שקולים:
\begin{enumerate}[label=(\alph*)]
\item 
$A$
שקולת שורה ל-
$I_n$.


\item 
$\rank(A)=n$


\item 
לממ"ל
$Ax=\underline{0}$
יש פתרון יחיד.


\item 
לכל
$b\in\F^n$
יש פתרון יחיד לממ"ל
$Ax=b$.


\item 
לכל
$b\in\F^n$
יש פתרון לממ"ל
$Ax=b$.


\end{enumerate}
\bigskip
\bigskip

\noindent\textbf{הגדרה 6.32.}

מטריצה 
$A\in\M_{n\times n}(\F)$
נקראת הפיכה משמאל אם קיימת
$B\in\M_{n\times n}(\F)$
כך ש-
$BA=I_n$.
$A$
נקראת הפיכה מימין אם קיימת
$C\in\M_{n\times n}(\F)$
כך ש-
$AC=I_n$.
אם שני התנאים מתקיימים, 
$A$
נקראת הפיכה.
\bigskip
\bigskip

\noindent\textbf{משפט 6.33.}

לכל
$A\in\M_{n\times n}(\F)$
התנאים הבאים שקולים:
\begin{enumerate}[label=(\alph*)]
\item 
$A$
שקולת שורה ל-
$I_n$.


\item 
$\rank(A)=n$


\item 
לממ"ל
$Ax=\underline{0}$
יש פתרון יחיד.


\item 
לכל
$b\in\F^n$
יש פתרון יחיד לממ"ל
$Ax=b$.


\item 
לכל
$b\in\F^n$
יש פתרון לממ"ל
$Ax=b$.


\item 
$A$
הפיכה משמאל.

\item 
$A$
הפיכה מימין.

\item 
$A$
הפיכה.


\end{enumerate}
\bigskip
\bigskip

\noindent\textbf{טענה 6.34.}

יהיו 
$A,B,C\in\M_{n\times n}(\F)$ 
כך ש-
$BA=I_n=AC$.
אז
$B=C$.
\bigskip
\bigskip

\noindent\textbf{מסקנה 6.35.}

יהיו 
$A,B\in\M_{n\times n}(\F)$.
אז מתקיים
$$
.AB=I_n \iff BA=I_n
$$
\bigskip
\bigskip

\noindent\textbf{מסקנה 6.36.}

יהיו
$A,B_1,B_2\in\M_{n\times n}(\F)$
כך ש-
$B_1A=B_2A=I_n$.
אז 
$B_1=B_2$.
\bigskip
\bigskip

\noindent\textbf{הגדרה 6.37.}

תהי
$A\in\M_{n\times n}(\F)$
הפיכה. נסמן ב-
$A^{-1}$
את המטריצה שמקיימת
$$
.AA^{-1}=A^{-1}A=I_n
$$
נקרא לה המטריצה ההופכית ל-
$A$.
\bigskip
\bigskip

\noindent\textbf{מסקנה 6.38.}

\label{cor:elementary-product}\
\begin{enumerate}[label=\alph*.]
\item
תהי 
$A\in\M_{n\times n}(\F)$
הפיכה. אז מתקיים
$$
.A^{-1}=E_kE_{k-1}\cdots E_1
$$
כאשר
$E_1,E_2,...,E_k$
הן המטריצות האלמנטריות המתאימות לפש"אות
$S_1,S_2,...,S_k$
בדירוג הקנוני של 
$A$.


\item 
$A$
היא מכפלה של מטריצות אלמנטריות:
$$
.A=E_1^{-1}E_2^{-1}\cdots E_k^{-1}
$$


\end{enumerate}
\bigskip
\bigskip

\end{document}